\chapter{MARL与传统规控的混合}
\label{ch:mr-marl-hybrid}

% ════════════════════════════════════════════════════════════════
%  全吸收章：重渲染 Tutorial 笔记
%   - 04_移动机器人规控/50_多机器人协作/130_MARL与传统规控混合.md
%     （混合范式动机/四层架构/RL→MPC 接口/CBF 安全过滤/残差 RL/Safe MARL/LLM 高层）
%  记号归一：R∈SO(3)、T∈SE(3)（preamble 宏 \SO \SE \so \se）；控制/RL/博弈自然记号保留。
%  去重：MARL 基础（CTDE/非平稳/HAPPO/MAPPO/值分解）已在 ch:plan-safemarl 深讲，本章一律
%        \cref{ch:plan-safemarl} 指过去，不重述；CBF/GNE/GCBF+ 的“安全=博弈均衡”原理同。
%        共识/ADMM 见 ch:mr-consensus；采样式 MPC/扩散 MPPI 见 mppi 章。
%  源由 AI 生成，论文归属/年份/会议/数值/前沿编号存疑处就地 \pz/\rebuilt，并入 claims 台账。
%  教学层遵循 v5.0 理论教学规范（动机→反面→历史→理论→陷阱→练习 + 误解汇总/小结/速查/故障排查）。
% ════════════════════════════════════════════════════════════════

\rebuilt{本章重渲染自 Tutorial 笔记《MARL 与传统规控混合》，其源稿由 AI 撰写，论文归属、年份、卷期与前沿工作（GCBF+、Residual MPC、CBF-RL、SHIELD、RoCo 等）的编号与实验数值可能有误。本章忠实重述、不擅自“修正”，逐处 \texttt{\textbackslash cite} 标源；存疑断言已登记 claims 台账，待联网核对后二次校验。}

\begin{practice}[前置自测]
开始之前，先试答下列问题。若已能流畅作答，可只速读本章表格与速查卡；若卡住（答不出 $\ge 2$ 题），标①者回 \cref{ch:plan-safemarl}（MARL 全栈：CTDE/非平稳/HAPPO/MAPPO/CBF-QP），标②者回控制部 MPC、标③者回 \cref{ch:plan-safemarl} 的 CBF 节，标④者回任务分配/MAPF 章，再来读本章。本章\textbf{大量复用}这些前置，但不从零重教——它教的是“如何把它们拼起来”。
\begin{enumerate}
  \item \textbf{（接 \cref{ch:plan-safemarl}）}① 什么是\textbf{集中训练分散执行}（CTDE）？为什么一个训练好的 MARL 策略 $\pi_\theta(a_i\mid o_i)$ 的输出，在真机上\textbf{不能}直接当电机力矩用？
  \item \textbf{（接控制部）}② 模型预测控制（MPC）的“滚动时域”是什么意思？MPC 的\textbf{硬约束}（关节限位、摩擦锥）与\textbf{软约束}（代价惩罚）在数学形式上有何本质区别？
  \item \textbf{（接 \cref{ch:plan-safemarl}）}③ 控制屏障函数（CBF）的核心不等式 $\dot h(x,u)\ge-\alpha(h(x))$ 保证什么？\textbf{相对度}是什么——当 $h$ 不显含 $u$ 时为何需要高阶 CBF（HOCBF）？
  \item \textbf{（接规划部）}④ 多机路径规划（MAPF）的“谁去做哪件事”这个离散决策，和“怎么走”这个连续控制，为什么要分成两层来解？
  \item \textbf{（综合）}假设你已有一个能在仿真里完美协调 $8$ 台无人机穿越障碍的 MARL 策略，现在要上真机。你最担心的\textbf{三件事}分别是什么？（提示：训练分布 vs 真实分布、安全保证、可解释性。）
\end{enumerate}
\noindent\emph{答案线索}：1$\to$\cref{sec:mh-why}（软决策的边界）；2$\to$\cref{sec:mh-hierarchy}（时间尺度分离）；3$\to$\cref{sec:mh-safety}（CBF 与相对度）；4$\to$\cref{sec:mh-why}（组合爆炸与分层）；5$\to$\cref{sec:mh-why}（三类失败 $+$ 四类 sim-to-real gap）。
\end{practice}

\section*{本章目标}
学完本章后，你应当能够：
\begin{enumerate}
  \item \textbf{论证}混合范式的必然性：从“探索-保证两难”出发，严格说清纯 MARL 与纯传统规控各自的能力边界，以及“分层解耦”为何是化解两难的工程主轴；
  \item \textbf{设计}分层混合架构：把一个多机器人任务分解到“高层决策（RL）$\to$ 中层 MPC $\to$ 底层 WBC/PD $\to$ 安全过滤层（CBF）”四个层次，说清每层的输入输出、时间尺度与归属理由；
  \item \textbf{实现} RL 高层 $+$ MPC 底层的接口：讲清 RL 策略输出（目标速度/编队偏移/子目标）如何变成 MPC 的参考/代价，并理解“接口语义对齐”为何是系统成败的关键；
  \item \textbf{推导}安全过滤器：从 CBF-QP 出发，推导“最小修正”的 KKT 闭式解，理解它就是\emph{到安全半空间的正交投影}，并理解 GCBF+ 如何用图神经网络把保证扩展到上千 agent；
  \item \textbf{掌握}残差强化学习：讲清 $u=u_{\text{base}}+u_{\text{res}}^{\theta}$ 为何样本效率高、为何继承基线的\emph{安全性下界}，并能量化残差幅值 $\rho$ 与基线裕度的关系；
  \item \textbf{公式化} Safe MARL：把安全需求写成约束马尔可夫博弈（CMG），用 Lagrangian 松弛 $+$ 原始-对偶求解，理解“期望约束”与“逐状态硬保证”两种安全语义的根本差异；
  \item \textbf{做出}有依据的范式选型：沿“是否需硬安全 $\times$ 是否有可靠模型 $\times$ 任务组合复杂度 $\times$ 规模”四维判断该用纯 MARL、纯规控还是某种具体混合形态，并理解高层决策器（MARL/LLM）的可替换性。
\end{enumerate}

\subsection*{本章知识导航}
\textbf{一句话定位}：本章把“多机协调”的两条技术路线——\cref{ch:plan-safemarl} 的\emph{学习路线}（MARL）与控制部/规划部的\emph{模型路线}（MPC/CBF/MAPF）——\textbf{焊接}成一个分层混合系统。核心主张只有一句话：\textbf{让 RL 在它最强的地方（高层、长程、难建模的协调决策）探索，让传统规控在它最强的地方（底层、实时、有模型的执行与安全）保证}，两者通过一个语义清晰的接口解耦。

\begin{center}
\small
\begin{tikzpicture}[
  node distance=5mm and 7mm, >=Stealth,
  blk/.style={draw, rounded corners, align=center, font=\small, inner sep=3pt, fill=main!6, draw=main!60!black},
  grp/.style={draw, rounded corners, align=center, font=\small, inner sep=3pt, fill=second!6, draw=second!60!black},
  fin/.style={draw, rounded corners, align=center, font=\small, inner sep=3pt, fill=third!8, draw=third!60!black},
  ln/.style={->, thick, gray!60!black}]
\node[blk] (why) {为什么混合\\ 探索-保证两难};
\node[blk, right=of why] (hier) {分层混合\\ 四层架构};
\node[grp, right=of hier] (iface) {RL$\to$MPC 接口\\ 命令空间};
\node[fin, below=9mm of why] (safe) {安全过滤\\ CBF-QP / GCBF+};
\node[fin, below=9mm of hier] (res) {残差 RL\\ $u_{\text{base}}+u_{\text{res}}$};
\node[blk, below=9mm of iface, fill=main!8, draw=main!60!black] (sel) {Safe MARL\\ $+$ 三范式选型};
\draw[ln] (why)--(hier); \draw[ln] (hier)--(iface);
\draw[ln] (why.south)--(safe.north);
\draw[ln] (hier.south)--(res.north);
\draw[ln] (iface.south)--(sel.north);
\draw[ln] (safe.east)--(res.west);
\draw[ln] (res.east)--(sel.west);
\end{tikzpicture}
\end{center}

\noindent\textbf{推荐路径}：\cref{sec:mh-why}（为什么混合）是\emph{思想地基}（必读，决定你是否真懂这一章在干什么）$\to$ \cref{sec:mh-hierarchy}（四层架构）给出\emph{骨架}（必读，后面所有内容都挂在这张图上）$\to$ \cref{sec:mh-interface}（RL$\to$MPC 接口）是接合点（做分层混合必读、需动手）$\to$ \cref{sec:mh-safety}（CBF 过滤 $+$ GCBF+）是\emph{安全硬核}（做安全关键系统必读且需跟着推一遍）$\to$ \cref{sec:mh-residual}（残差 RL）是另一条混合路线（做精细操作/腿足必读）$\to$ \cref{sec:mh-safemarl}（Safe MARL）把安全写进训练目标并收口选型 $\to$ \cref{sec:mh-llm}（LLM 高层）是前沿延伸。\cref{sec:mh-safety} 与 \cref{sec:mh-residual} 是两个最硬核小节，值得花最多时间。

\subsection*{前置知识桥接}
本章紧接 \cref{ch:plan-safemarl}（学习路线）与控制部/规划部（模型路线），这里把最关键的前置点重新激活——你不必翻回去就能跟上：
\begin{itemize}
  \item \textbf{MARL 输出的是“软”决策}（\cref{ch:plan-safemarl}）。那里用 CTDE 降低非平稳性、用 HAPPO 处理异构智能体的序贯单调更新。要记住的是：MARL 学到的 $\pi_\theta(a_i\mid o_i)$ 是一个\textbf{概率策略}，它通过奖励塑形把“碰撞”惩罚成低概率事件，但\textbf{永远无法把碰撞概率压到零}——奖励是软的。本章要给这个“软决策”套上一层“硬保证”的外壳。
  \item \textbf{MPC 是带模型的实时优化器}（控制部）。MPC 的力量来自\textbf{它知道动力学模型}——它能在预测窗口内“看到未来 $H$ 步”，并在硬约束（摩擦锥、关节限位）下求出物理可行的控制。本章让 MARL 不再直接碰这些底层约束，而是把它们留给 MPC。
  \item \textbf{安全 $=$ 集合的前向不变性}（\cref{ch:plan-safemarl} 的 CBF 节）。CBF 把“安全”形式化为“状态永远留在 $\mathcal{C}=\{x:h(x)\ge0\}$ 内”，并给出一个\emph{逐时刻}的控制约束 $\dot h\ge-\alpha(h)$ 来保证它——这是一个比“MPC 软约束”和“MARL 奖励惩罚”都更强的、\emph{逐状态可证明}的保证。\cref{sec:mh-safety} 的安全过滤层就建立在它之上。
  \item \textbf{分层是对抗组合爆炸的标准手段}（规划部）。任务分配与 MAPF 已经把“谁去做”与“怎么走”分成两层，正因离散决策与连续控制混在一起会爆炸。本章把同样的“分层解耦”思想推广到“学习 vs 模型”这个新的切分维度。
\end{itemize}

\textbf{前向预告}：本章的混合架构是整个多机器人方向的\textbf{收口}——它把前面所有章节（通信图、共识、任务分配、MAPF、编队、分布式 MPC、MARL、CBF）的成果组织成一个可部署的系统。读完本章，你应能回答一个总览性问题：\textbf{“给我一个真实的多机器人任务，我该用前面学的哪些模块、怎么拼？”} \cref{sec:mh-safemarl} 的选型框架会正面回答它。

\subsection*{如果跳过本章会怎样}
\begin{itemize}
  \item \textbf{仿真里完美，一上真机就撞}：你用 \cref{ch:plan-safemarl} 的 MARL 训出一个能在仿真里漂亮协调 $6$ 台无人机的策略，部署到真机——结果某机在训练时从没出现过的风扰下输出一个剧烈动作，直接撞向邻机，而你\textbf{没有任何机制能阻止}这次碰撞（奖励里的碰撞惩罚只在训练时起作用）。\cref{sec:mh-why} 会告诉你为什么“纯 MARL 上真机”是范式层面的错误，\cref{sec:mh-safety} 会告诉你怎么用 CBF 安全过滤层从根本上杜绝这类碰撞。
  \item \textbf{纯 MPC 写不出协调策略，调参调到崩溃}：你反过来想全用 MPC，给 $8$ 台异构机器人写集中式分布式 MPC 协同搬运并避障，很快发现“谁让谁、编队怎么变形”这种\textbf{高层、组合、难解析建模}的决策根本写不进代价函数——你用一堆手工权重去逼近它，结果权重互相打架，调一个崩一个。\cref{sec:mh-why} 会告诉你为什么这类决策恰是 MARL 的主场，\cref{sec:mh-interface} 会告诉你怎么让 MARL 出高层命令、MPC 只负责执行。
\end{itemize}

\begin{table}[H]\centering\small
\caption{本章阅读方式建议}
\label{tab:mh-reading}
\begin{tabular}{lll}
\toprule
阅读方式 & 时间 & 适合谁 \\
\midrule
精读（推完 CBF-QP 的 KKT 解、搭分层接口、实现 pairwise CBF 过滤 $+$ 残差叠加、推完 CMG 的 Lagrangian） & 18--24 小时 & 第一次系统学混合范式 \\
速读（读懂四层架构、三种混合形态的适用边界、关键公式，跳过推导细节） & 5--7 小时 & 已有 MARL $+$ MPC 背景 \\
速查（架构图 $+$ 安全过滤 KKT 解 $+$ 残差公式 $+$ 三范式对比 $+$ 选型框架 $+$ 故障排查） & 50--70 分钟 & 设计混合系统时回查 \\
\bottomrule
\end{tabular}
\end{table}

\begin{note}
\textbf{本章记号与约定.}\;
旋转矩阵记 $\mathbf{R}\in\SO(3)$、位姿 $\mathbf{T}\in\SE(3)$（preamble 宏 $\SO,\SE,\so,\se$，与全书一致，\nameref{ch:notation}）；控制/强化学习领域的自然记号保留。
智能体数 $N$，下标 $i,j$；第 $i$ 个 agent 的局部观测 $o_i$、动作 $a_i$、状态 $x_i$、控制 $u_i$、位置 $p_i$；MARL 策略 $\pi_\theta(a_i\mid o_i)$。
控制仿射动力学 $\dot x=f(x)+g(x)u$（$f$ 漂移项、$g$ 输入项）。
安全集 $\mathcal{C}=\{x:h(x)\ge0\}$，CBF 标量函数 $h$，扩展 class-$\mathcal{K}$ 函数 $\alpha$；李导数 $L_fh,L_gh$。
\textbf{字母隔离声明}：本章 $\alpha$ 指 CBF 的 class-$\mathcal{K}$ 系数（非学习率，学习率记 $\eta$）、$\mu$ 指 Lagrangian 乘子、$\rho$ 指残差幅值上限、$\delta$ 指累积代价的安全阈值、$\gamma$ 为折扣因子、$H$ 为 MPC 预测时域、$c_{i,t}$ 为安全代价，均仅本章局部生效。
\end{note}

\begin{note}
\textbf{科研发展脉络（一图胜千言）.}\;\rebuilt{下表年份/会议/数值转述自源稿，待联网核对，勿张冠李戴。}
混合范式不是某篇论文的发明，而是两条独立路线（学习 vs 模型）在“都撞到自己天花板”后被迫融合的结果，有三股力量在汇流：
\textbf{(一) 安全过滤一侧}——Ames 的 CBF（综述 ECC 2019\cite{ames2019cbf}）把安全做成可\emph{叠加在任何控制器之上}的 QP 模块 $\to$ GCBF+（T-RO 2024\cite{zhang2025gcbfplus}）用 GNN 扩展到上千 agent $\to$ CBF-RL / SHIELD（2025\cite{gandhi2023shield,yin2023shield}）把它从“部署期外挂”推进到“训练期内化”，主线是\emph{安全的可扩展性与与学习的融合度不断提升}；
\textbf{(二) 残差学习一侧}——Johannink 的 residual RL（ICRA 2019\cite{johannink2019residual}）确立“在可信基线上只学修正”$\to$ Residual MPC（2025\cite{residualmpc2025}）把基线从简单 PD 升级为 GPU 并行 MPC，主线是\emph{基线越来越强，RL 只补“模型补不了的那一点”}；
\textbf{(三) 分层混合一侧}——高层 RL $+$ 底层控制的雏形 % \cite{silver2005cooperative} 误引：该键指向 David Silver《Cooperative Pathfinding》(AIIDE 2005) 的多机寻路论文，与“高层 RL+底层控制”无关，疑为与 Tom Silver《Residual Policy Learning》(2018) 混淆。\rebuilt{此处“分层混合雏形”引用已注释，待补正确出处，勿张冠李戴。}
$\to$ MPC $\times$ predictive RL（2023\cite{bjelonic2023mpcrl}，RL 出尾代价）$\to$ RoCo（ICRA 2024\cite{mandi2024roco}，LLM 出高层决策），主线是\emph{高层决策器越来越“聪明”、底层执行器越来越“可靠”}。
权威综述见 Brunke 等\emph{Safe Learning in Robotics}\cite{brunke2022safe}。三股力量的共识：\textbf{没人再相信“纯学习”或“纯模型”能单独解决真实多机器人问题}，分歧只在“在哪个层次、用什么方式把两者缝合”。
\end{note}

% ════════════════════════════════════════════════════════════════
\section{为什么要混合：探索-保证两难}
\label{sec:mh-why}

\paragraph{动机：一个无法两全的真实困境。}
设想一个具体任务：\textbf{$8$ 台异构机器人（$4$ 架无人机 $+$ $4$ 台地面车）要在一个有动态障碍的仓库里协同把货物从 A 区搬到 B 区}。这个任务里同时存在两类截然不同的子问题。\textbf{第一类是“高层协调决策”}：哪台车配合哪架无人机？遇到只能容一台车通过的狭窄通道谁先谁后？某架无人机电量低了编队怎么动态重组？这些决策\emph{长程}（好坏要很多步后才看得出）、\emph{组合性强}（选择空间随机器人数与任务数爆炸）、\emph{难以解析建模}（“什么样的让行顺序最好”没有干净的数学表达式）。\textbf{第二类是“底层执行控制”}：每台机器人怎么精确跟踪目标轨迹？怎么在摩擦锥、关节限位、力矩上限内算出可行控制？怎么保证两台机器人物理上\emph{绝不相撞}？这些问题特征完全相反——\emph{短程}（逐周期反馈）、\emph{连续}（实数控制空间）、\emph{有精确模型}（动力学/碰撞几何都可写出方程）。问题来了：\textbf{你能用同一种方法把这两类子问题都解好吗？} 答案是不能——这正是混合范式存在的根本原因。

\paragraph{反面一：如果用纯 MARL 会怎样。}
纯 MARL 把整个任务（从高层协调到底层力矩）交给一个端到端多智能体策略 $\pi_\theta$，用精心设计的奖励训练它。这条路在仿真里常常惊艳，但一旦面对真机，会在三处系统性失败。
\textbf{失败点一：没有任何硬安全保证（最致命）。} MARL 的安全通过\emph{奖励塑形}实现——在奖励里加“碰撞 $=-1000$”。但奖励是“软”的：它让策略\emph{倾向于}避免碰撞、把碰撞变成低概率事件，却\textbf{永远无法把碰撞概率压到零}。仿真里 $99.9\%$ 不碰撞听起来很好，但真机上 $0.1\%$ 的碰撞概率意味着平均每 $1000$ 次操作就损毁一次硬件——工业部署完全不可接受。
\textbf{失败点二：分布偏移下的灾难性失效。} MARL 策略在仿真动力学下训练，真实世界存在多类系统性分布偏移，任一类都可能把策略推到训练时从未见过的\emph{分布外}（out-of-distribution, OOD）观测；黑盒网络在 OOD 输入上的行为\textbf{完全不可预测}，可能输出剧烈危险动作，且多智能体里一个 agent 的异常会通过交互\emph{传染}给其他 agent。
\textbf{失败点三：不可解释、难调试、无法认证。} 真机出问题时你面对一个数百万参数的黑盒，没有任何可检查的中间量（没有“它认为的目标轨迹”、没有“它估计的障碍位置”），调试变成玄学；安全关键领域还意味着\textbf{无法通过认证}。

\begin{insight}[奖励惩罚与 CBF 约束之间有一道范式鸿沟]
MARL 的奖励惩罚是\textbf{期望意义}的优化目标（“平均而言少碰撞”），CBF 约束是\textbf{逐状态}的硬保证（“这一帧绝不越界”）。你无法通过把碰撞惩罚调到无穷大把前者变成后者——惩罚再大，只要它是奖励的一部分，策略就可能为了某个更大的累积奖励而“赌一把”小概率碰撞，这是优化的数学本质决定的，不是调参能解决的。\textbf{安全不能是优化目标的一项，它必须是优化的约束。}
\end{insight}

\paragraph{反面二：如果用纯传统规控会怎样。}
反过来全用传统规控（MPC $+$ CBF $+$ MAPF $+$ 手工协调规则）行不行？传统规控恰好补上 MARL 缺的三块——有硬安全保证、模型准确时行为可预测、每个中间量都可解释。但它在另一处系统性失败：\textbf{无法处理“难以解析建模”的高层协调决策}。
\textbf{失败点一：高层协调写不进代价函数。} 传统规控的核心是“把问题写成优化问题然后求解”，但“$8$ 台异构机器人遇狭窄通道时的最优让行策略”\emph{没有干净的数学表达式}。你只能用两种笨办法逼近：要么堆一大堆 \verb@if-else@ 启发式规则（情境一复杂就漏洞百出），要么往 MPC 代价里堆手工权重项（陷入“权重打架”——调高“编队保持”，“避障灵活性”就崩，反之亦然）。
\textbf{失败点二：组合爆炸。} $N$ 个机器人的联合规划随 $N$ 指数爆炸，MAPF 的最优求解是 NP-难的；近似策略本身也难以手工设计（“什么样的优先级排序最好”又回到“难以解析建模”）。
\textbf{失败点三：对未建模动态束手无策。} MPC 的力量来自模型，但\emph{模型之外的东西它一概看不见}——复杂摩擦、柔性形变、空气动力学扰动、其他机器人的真实（而非假设的）行为……这些“残差动力学”无法精确建模，纯 MPC 只能当噪声忽略，导致跟踪精度下降。

\begin{insight}[“写成优化问题”既是规控的力量也是它的天花板]
传统规控的力量来自“把问题写成优化问题”，但这也正是它的天花板——\textbf{有些决策本质上无法被简洁地写成优化问题}。最优让行、最优编队重组、最优任务再分配，其“最优”依赖海量难以枚举的情境，价值函数高度非线性、非凸、甚至不连续。强行用手工代价逼近，本质是在用人脑做策略网络该做的事，而人脑做不好这件事。这恰是“学习”的主场：RL 能从经验中\emph{自动发现}这个复杂价值函数的形状。
\end{insight}

\subsection{理论：解耦原理与 sim-to-real gap 的局部化}
\label{sec:mh-decoupling}

现在把上面的“故事”提炼成可操作的理论。混合范式的核心理论支柱只有一条——\textbf{解耦原理}（decoupling principle）：把一个复杂系统沿“特征维度”切分成若干层，让每一层只处理它有比较优势的子问题，层与层之间通过一个\emph{窄而清晰的接口}通信。关键是识别出“学习”和“模型”各自的比较优势维度，它们恰好在多个维度上互补：

\begin{table}[H]\centering\small
\caption{学习（MARL）与模型（规控）的比较优势维度}
\label{tab:mh-decouple-dims}
\begin{tabular}{lll}
\toprule
维度 & 适合学习（MARL）的一端 & 适合模型（规控）的一端 \\
\midrule
时间尺度 & 慢（$\sim 10$ Hz，决策不需高频） & 快（$50$--$1000$ Hz，控制需实时反馈） \\
抽象层次 & 高层（编队、让行、任务分配） & 底层（力矩、轨迹跟踪） \\
建模难度 & 难建模（协调价值函数无解析式） & 易建模（动力学/碰撞有方程） \\
决策类型 & 探索性（需从经验发现策略） & 保证性（需可证明的安全/稳定） \\
空间结构 & 组合/离散（谁配合谁、谁先过） & 连续（实数控制空间优化） \\
\bottomrule
\end{tabular}
\end{table}

\noindent\textbf{解耦的工程含义}：把“难建模、高层、慢、探索性、组合”的部分交给 MARL，把“易建模、底层、快、保证性、连续”的部分交给规控，两者通过一个\emph{语义清晰的命令接口}（如“目标质心速度”、“编队偏移”、“子目标位置”）通信。这个接口必须满足：MARL 输出的命令，无论多离谱，底层规控都能\textbf{安全地}尝试执行（哪怕执行不到，也不会违反硬约束）——这个“安全可执行”的性质正是 \cref{sec:mh-safety} 安全过滤层要保证的。

\begin{insight}[解耦的深层价值：它把 sim-to-real gap 局部化了]
在纯 MARL 里，sim-to-real gap 弥漫在从感知到力矩的整条链路上，无处不在、无法隔离。而在分层混合里，MARL 只负责高层决策（$\sim 10$ Hz 的编队/让行命令），底层执行用基于真实模型的 MPC。于是 \textbf{sim-to-real gap 被压缩到“高层决策”这一个窄接口上}——而高层决策对底层动力学细节天然不敏感（“让某机靠左一点”这个决策，在仿真和真机里含义几乎一样）。这就是为什么混合架构的迁移比纯 MARL 容易得多：你只需让一个低维、抽象、对动力学不敏感的决策迁移，而不是让整条感知-控制链路迁移。
\end{insight}

为了让“分布偏移”这个抽象概念变得可操作，对 sim-to-real gap 做\emph{穷举式的维度分类}（一个“思考框架”而非“记忆清单”）：

\begin{table}[H]\centering\small
\caption{四类 sim-to-real gap 及混合架构的化解层}
\label{tab:mh-gap}
\begin{tabular}{p{2.0cm}p{3.6cm}p{3.2cm}p{3.4cm}}
\toprule
维度 & 具体来源 & 对纯 MARL 的威胁 & 混合架构如何缓解 \\
\midrule
动力学 gap & 质量/惯量/摩擦估计误差、接触模型简化 & 策略在错误动力学下学的动作在真机失效 & 底层 MPC 用真实标定模型，gap 不影响高层决策 \\
感知 gap & 传感器噪声、延迟、遮挡、点云稀疏 & OOD 观测让策略输出危险动作 & 高层决策对感知细节不敏感；安全过滤用实时局部感知 \\
执行 gap & 电机延迟、量化误差、力矩饱和 & 策略假设的理想执行在真机打折 & 底层 WBC/PD 在真实执行约束内求解 \\
多体交互 gap & 其他 agent 的真实行为 $\ne$ 训练时假设行为 & 一个 agent 异常通过交互传染全队 & CBF 过滤保证 pairwise 安全，切断传染链 \\
\bottomrule
\end{tabular}
\end{table}

\noindent 理解这张表的关键是：\textbf{纯 MARL 必须独自承受全部四类 gap，而混合架构把每一类 gap 都甩给了最适合处理它的层}。MARL 只需处理它最擅长的“高层协调决策”，而这一层对四类 gap 都相对鲁棒。

\subsection{三种混合形态的预览}
\label{sec:mh-three-forms}

基于解耦原理，混合范式在工程上结晶出三种主要形态。这里先给全景预览，后续小节逐一展开——你现在只需记住“有这三种，它们解耦的方式不同”：

\begin{table}[H]\centering\small
\caption{三种混合形态：同一解耦原理沿三个维度的切分}
\label{tab:mh-forms}
\begin{tabular}{p{2.4cm}p{3.4cm}p{2.8cm}p{2.0cm}p{2.4cm}}
\toprule
形态 & 核心思想 & 解耦方式 & 本章小节 & 典型代表 \\
\midrule
分层混合 & RL 出高层命令，规控管底层执行 & 按抽象层次纵向切分 & \cref{sec:mh-hierarchy,sec:mh-interface} & RL$+$MPC 四足、RoCo \\
安全过滤 & RL 出名义动作，CBF 过滤成安全动作 & 按“决策 vs 安全”切分 & \cref{sec:mh-safety} & GCBF+、CBF-RL、SHIELD \\
残差学习 & 控制器出基线，RL 只学修正量 & 按“基础 vs 残差”叠加 & \cref{sec:mh-residual} & residual RL、Residual MPC \\
\bottomrule
\end{tabular}
\end{table}

\noindent 这三种形态不是互斥的——一个成熟系统往往\emph{同时用上多种}（分层混合的高层是 RL、底层是 MPC，同时在 MPC 输出上叠一个 CBF 过滤，甚至 MPC 本身还带一个 residual RL 修正）。\cref{sec:mh-safemarl} 会教你如何根据任务特征组合它们。

\begin{insight}[分层、过滤、残差是同一原理的三种刀法]
“分层、过滤、残差”本质上是同一个解耦原理沿三个不同维度的切分：分层沿\textbf{抽象层次}（上/下）切，过滤沿\textbf{职能}（决策/安全）切，残差沿\textbf{贡献}（基础/修正）切。它们都在回答同一个问题——“让 RL 做什么、让模型做什么、两者怎么对接”——只是切分的刀法不同。理解了这一点，你就不会把它们当成三个孤立的技巧，而会看到背后统一的设计哲学。
\end{insight}

\begin{pitfall}[概念误区：以为“把碰撞惩罚调大”就能让 MARL 安全]
错误做法：把碰撞惩罚设成 $-10^4$ 甚至 $-10^9$，认为策略就“绝对不会”碰撞。后果：训练时碰撞率下降但永远不为零；部署仍偶发碰撞，且惩罚过大还导致策略过度保守（为躲远处障碍而完全不动）或训练不稳定（梯度爆炸）。根本原因：奖励是\emph{期望意义}的软目标，碰撞是它优化的一个\emph{项}而非\emph{约束}。正确做法：把安全从“奖励项”提升为“硬约束”——用 CBF 安全过滤层（\cref{sec:mh-safety}）在动作执行前强制投影到安全集，或用约束马尔可夫博弈（\cref{sec:mh-safemarl}）把安全写成显式约束 $\mathbb{E}[\text{cost}]\le\delta$。安全和性能必须在\textbf{不同的数学层次}上处理。
\end{pitfall}

\begin{pitfall}[思维陷阱：以为“纯 MPC 更安全所以应该尽量多用 MPC”]
错误做法：因为 MPC 有安全保证，就把高层协调决策也塞进 MPC，用手工代价 $+$ 启发式逼近。后果：代价权重互相打架，调一个崩一个；启发式在复杂情境下漏洞百出。根本原因：混淆了“安全”和“决策质量”两个正交维度。MPC 在“有模型的执行”上确实更安全，但在“难建模的协调决策”上\emph{根本没有优势}——它只是把“学习策略”换成“人脑手工设计代价”。正确做法：让 MPC 只做它有比较优势的事（底层执行 $+$ 硬约束），把高层协调交给 RL，“安全”由专门的安全层（CBF）保证，与“决策”解耦。
\end{pitfall}

\begin{pitfall}[思维陷阱：把“混合架构”理解成“先跑 MARL 再跑 MPC”的简单串联]
错误做法：以为混合就是“MARL 算一个结果喂给 MPC 再算一遍”的流水线串联，不考虑接口语义与时间尺度。后果：接口语义不对齐（MARL 输出的“动作”MPC 不知道怎么用），或时间尺度错配（MARL $10$ Hz 的命令被 MPC $50$ Hz 反复执行导致抖动），系统行为混乱。根本原因：没理解解耦原理的核心是“沿正确维度切分 $+$ 设计语义清晰的窄接口”。正确做法：先想清解耦维度，再设计接口语义，最后处理频率匹配（\cref{sec:mh-hierarchy,sec:mh-interface}）。
\end{pitfall}

\begin{practice}
\begin{enumerate}
  \item \textbf{【概念辨析】} 用你自己的话解释“为什么把碰撞惩罚调到无穷大也不能让 MARL 真正安全”，要求：(a) 从“优化目标项 vs 优化约束”的角度论证；(b) 举一个具体反例情境，说明策略可能为更大奖励而“赌”一次小概率碰撞；(c) 说明 CBF 安全约束在数学形式上和“无穷大惩罚”的本质区别。
  \item \textbf{【设计分析】} 回到“$8$ 台异构机器人仓库搬运”任务，把所有子问题\emph{显式列出}（至少 $6$ 个），为每个判断它该交给 MARL 还是传统规控，并沿 \cref{tab:mh-decouple-dims} 的五维分析判断依据。
  \item \textbf{【反事实推理】} 本节论证了“分层混合把 sim-to-real gap 局部化到高层决策接口上”。请做反事实分析：若把切分点放错——让 MARL 直接输出\emph{关节力矩}、底层只做简单力矩限幅——这个“混合”还能享受 gap 局部化的好处吗？为什么？（提示：力矩这个输出对动力学 gap 的敏感度。）
\end{enumerate}
\end{practice}

% ════════════════════════════════════════════════════════════════
\section{分层混合架构：四层分工}
\label{sec:mh-hierarchy}

\paragraph{动机：从“两类子问题”到“多层流水线”。}
\cref{sec:mh-why} 把任务粗分成“高层协调决策”和“底层执行控制”两类。但真实系统里，从“编队怎么变”到“电机出多少力矩”之间横跨好几个抽象层次和时间尺度——直接从 $10$ Hz 的“编队偏移命令”跳到 $1000$ Hz 的“关节力矩”，中间缺了好几环。想象一个指令链：MARL 决定“某机应向左偏移 $0.5$ 米以避让”（一个\emph{抽象的高层意图}），但电机只认识“力矩 $=0.37$ N$\cdot$m”。从前者到后者需逐层翻译：(1)“向左偏移 $0.5$ 米”$\to$ 一条\emph{平滑可行的目标轨迹}；(2) 目标轨迹 $\to$ 一串\emph{满足动力学和约束的控制参考}；(3) 控制参考 $\to$ 具体的\emph{关节力矩}；(4) 在第 2、3 步之间，还要有一道\emph{安全检查}拦下会导致碰撞的命令。这四步对应四个层次，每层有自己最合适的方法与频率。

\paragraph{历史：分层控制的思想源流。}
分层控制不是机器人学的新发明，其源头可追溯到 1970--80 年代的\emph{分层智能控制}（Saridis 的“组织-协调-执行”三级理论）和经典的\emph{串级控制}（cascade control，外环慢、内环快）。这些经典思想确立了一条原则：\textbf{用时间尺度分离来解耦一个复杂的多频率系统}——慢的外环管“目标”，快的内环管“跟踪”，只要内环比外环快足够多（一般 $5$--$10$ 倍），就可近似认为内环“瞬间”达到外环给的目标，从而把耦合系统拆成两个独立设计的子系统。机器人学把这个思想发扬光大：经典足式控制就是“步态规划（慢）$\to$ MPC（中）$\to$ 全身控制 WBC（快）”三层结构。混合范式做的事，是在这个成熟骨架的\textbf{最顶层}插入 RL——用学习替换掉原来最难手工设计的那一层（高层决策/协调），保留下面成熟可靠的模型层。

\begin{insight}[混合架构是“外科手术式”地把 RL 植入恰当位置]
分层混合最聪明的地方不是“发明了分层”，而是\textbf{精准识别出整个控制栈里“唯一一层最适合用学习替换”的位置}——最顶层的高层决策。底层（力矩、轨迹跟踪）有成熟的模型方法且要求实时和安全，不该动；唯有顶层的协调决策既难建模又不要求高频，正好是学习的甜区。混合架构是把 RL 植入成熟控制栈的恰当位置，而不是推倒重来。
\end{insight}

\subsection{理论：四层架构详解}
\label{sec:mh-four-levels}

下面给出多机器人 MARL $+$ 规控混合的标准四层架构（\cref{fig:mh-stack}），它是本章的骨架。

\begin{figure}[H]
\centering
\begin{tikzpicture}[
  >=Stealth, font=\small,
  lvl/.style={draw, rounded corners, align=left, inner sep=5pt, text width=10.4cm},
  l3/.style={lvl, fill=second!10, draw=second!60!black},
  l2/.style={lvl, fill=main!8, draw=main!60!black},
  lc/.style={lvl, fill=third!12, draw=third!60!black, text width=6.2cm},
  l1/.style={lvl, fill=main!6, draw=main!50!black},
  ln/.style={->, thick, gray!55!black}]
\node[l3] (l3) {\textbf{Level 3：MARL 高层决策}\hfill$\sim 10$ Hz\\[2pt]
  \footnotesize 输入：各 agent 局部观测 $o_i$（自身 $+$ 邻居 $+$ 任务）；输出：高层命令 $u_{\text{cmd}}$（质心目标速度 $v_{\text{cmd}}$ / 编队偏移 $\Delta p$ / 子目标）。方法：MAPPO/HAPPO（\cref{ch:plan-safemarl}）。为何在此：决策长程、组合、难建模 $\to$ 学习甜区；$10$ Hz 足够。};
\node[l2, below=7mm of l3] (l2) {\textbf{Level 2：单体 MPC（轨迹优化）}\hfill$\sim 50$ Hz\\[2pt]
  \footnotesize 输入：$u_{\text{cmd}}$ $+$ 当前状态 $x$；输出：力参考 / 关节轨迹 / 期望加速度。方法：滚动时域 OCP。为何在此：有精确动力学模型，能在硬约束（摩擦锥、关节限位）内求物理可行控制；$50$ Hz 平衡实时性与计算量。};
\node[lc, below left=9mm and -1mm of l2.south] (lc) {\textbf{安全过滤层：CBF-QP}\hfill\footnotesize 同频或更高\\[2pt]
  \footnotesize 输入：$u_{\text{mpc}}$ $+$ 邻居状态；输出：安全修正 $u_{\text{safe}}$。方法：CBF-QP / GCBF+。为何：逐状态硬安全保证（\cref{sec:mh-safety}）。\emph{横切层，非纵向一环}。};
\node[l1, below=22mm of l2] (l1) {\textbf{Level 1：WBC $+$ PD（全身控制）}\hfill$\sim 200$--$1000$ Hz\\[2pt]
  \footnotesize 输入：安全修正后的力参考 $u_{\text{safe}}$；输出：关节力矩 $\boldsymbol\tau$。方法：QP-WBC（零空间投影分层任务）$+$ 关节 PD。为何在此：把“力参考”翻译成满足整机动力学的关节力矩，需最高频率保证执行精度与稳定性。};
\node[draw, rounded corners, fill=gray!12, below=6mm of l1, text width=3cm, align=center] (body) {机器人本体};
\draw[ln] (l3) -- node[right, font=\scriptsize]{$u_{\text{cmd}}$（语义命令）} (l2);
\draw[ln] (l2.south) -- ++(0,-3mm) -| (lc.north);
\draw[ln] (lc.south) |- ($(l1.north)+(0,1mm)$);
\draw[ln] (l2.south) -- node[right, font=\scriptsize]{$u_{\text{mpc}}\!\to\!u_{\text{safe}}$} (l1);
\draw[ln] (l1) -- node[right, font=\scriptsize]{$\boldsymbol\tau$} (body);
\end{tikzpicture}
\caption{多机器人 MARL $+$ 规控混合的四层架构。频率阶梯（$10/50/200$--$1000$ Hz）由时间尺度分离原理决定；CBF 安全过滤层是\emph{横切}所有控制路径的守门员，而非流水线上的一个工位。}
\label{fig:mh-stack}
\end{figure}

\noindent\textbf{逐层解读}——为什么每层是现在这个频率、用现在这个方法：
\begin{itemize}
  \item \textbf{Level 3（MARL，$\sim 10$ Hz）}：混合架构里\emph{唯一的学习层}。为什么 $10$ Hz 就够？因为“编队怎么变、谁让谁”这类决策的\emph{时间常数本来就慢}（人开车做“变道决策”也就每秒一两次）。低频还带来三个好处：推理算力压力小（可用更大网络）、sim-to-real gap 被压缩在低维慢变接口上、一个决策能维持多个底层周期使行为更连贯。
  \item \textbf{Level 2（MPC，$\sim 50$ Hz）}：\emph{主执行层}。它拿到高层命令（如 $v_{\text{cmd}}$），结合精确动力学模型在预测时域 $H$ 内求解一个 OCP，输出物理可行的控制参考。为什么是 MPC？因为它能\emph{显式处理硬约束}（摩擦锥、关节限位）并\emph{预见未来}（在 $H$ 步窗口内提前规避）。$50$ Hz 是实时性与计算量的折中。
  \item \textbf{安全过滤层（CBF-QP）}：\emph{安全守门员}。它叠加在 MPC 输出之上，对 $u_{\text{mpc}}$ 做最小修正使其满足 CBF 安全约束。注意它是\textbf{横切层}而非纵向一环——本质是“控制下发到电机之前的最后一道硬安全检查”（\cref{sec:mh-safety} 详解）。
  \item \textbf{Level 1（WBC $+$ PD，$\sim 200$--$1000$ Hz）}：\emph{底层执行层}。它用全身控制（whole-body control, WBC，用零空间投影把多个任务分层）把“力参考”翻译成关节力矩，再经关节 PD 输出。为什么最高频？因为执行精度和闭环稳定性对延迟最敏感，底层跟踪环必须足够快才能压住高频扰动。
\end{itemize}

\begin{proposition}[时间尺度分离对频率阶梯的约束]\label{prop:mh-timescale}
若相邻两层频率比保持在 $5$--$10$ 倍（$10\to 50\to 200/1000$），则外环可把内环近似当作“理想的瞬时执行器”来设计，相邻两层从而\emph{近似解耦}、可独立设计；若两层频率太接近（如都是 $30$ Hz），时间尺度分离假设失效，两层会互相“看见”对方的暂态而产生耦合振荡。
\end{proposition}

\begin{insight}[频率阶梯不是拍脑袋定的]
四层架构的频率阶梯（$10/50/200$--$1000$ Hz）是\textbf{时间尺度分离原理的直接体现}（\cref{prop:mh-timescale}）。每相邻两层差一个数量级，保证上层可把下层近似当成“理想的瞬时执行器”来设计——这正是“分层”之所以能“解耦”的数学根基。如果你在自己的系统里发现两层频率太接近导致振荡，第一反应应是“时间尺度分离假设被破坏了”，而不是去调 PID 参数。
\end{insight}

\subsection{理论-工程桥接：四层对应到真实代码栈，与层间降级}
\label{sec:mh-stack-impl}

这套四层架构在真实开源栈里有清晰对应（\cref{tab:mh-impl}），理解这个对应能帮你把抽象图落地。

\begin{table}[H]\centering\small
\caption{四层架构在开源栈里的对应物（\rebuilt{库名/接口转述自源稿，待核}）}
\label{tab:mh-impl}
\begin{tabular}{p{1.8cm}p{2.6cm}p{4.2cm}p{3.2cm}}
\toprule
层次 & 抽象职能 & 典型开源实现 & 接口数据 \\
\midrule
Level 3 & MARL 高层决策 & 自训策略（MAPPO/HAPPO，基于 HARL 库） & $v_{\text{cmd}},\Delta p$ \\
Level 2 & 单体 MPC & OCS2（SQP-RTI）/ acados & 足端力 / 关节轨迹 \\
安全层 & CBF 过滤 & 自实现 CBF-QP（OSQP）/ GCBF+（JAX） & 修正后控制 $u_{\text{safe}}$ \\
Level 1 & WBC $+$ PD & OCS2 的 WBC 模块 / 自实现 QP-WBC & 关节力矩 $\boldsymbol\tau$ \\
\bottomrule
\end{tabular}
\end{table}

\noindent\textbf{关键工程点}：\emph{层与层之间用消息队列解耦，各层跑在独立的实时线程/进程里}，按各自频率运行。上层命令以“最新值”被下层采样（下层频率高，会重复用同一个上层命令多次，直到上层更新）。这种“异步多频率”结构正是时间尺度分离在软件工程上的落地——它让你可以独立开发、测试、替换每一层。

\textbf{一个常被忽略的设计：层间“降级安全”}。成熟的混合架构\emph{每一层都要能在“上层失效”时优雅降级}。如果 MARL 推理超时没给出新命令，Level 2 应保持上一个有效命令（或退回一个安全默认命令，如“原地悬停/停车”），而不是用过期或空命令；同理 MPC 求解失败（infeasible）时，Level 1 应有安全回退（如重力补偿 $+$ 阻尼）。这个设计的价值在于：\textbf{它把单点故障隔离在层内}。上层（学习层，最不可靠）的失效不会直接传导成物理事故，而是被下层（模型层，更可靠）的降级逻辑接住——这也是混合架构相比纯 MARL 的一个隐性安全优势（纯 MARL 没有“下层”可以降级，策略失效就是系统失效）。

\begin{pitfall}[思维陷阱：让相邻两层频率太接近，破坏时间尺度分离]
错误做法：为“让系统响应更快”，把 MARL 高层提到 $50$ Hz 和 MPC 同频。后果：高层和中层互相“看见”对方的暂态，产生耦合振荡——机器人出现莫名其妙的周期性抖动，且调任何单层参数都按下葫芦浮起瓢。根本原因：分层解耦依赖时间尺度分离假设（\cref{prop:mh-timescale}）。正确做法：保持相邻层 $5$--$10$ 倍频率比；如确需高层更快响应，应优化高层决策的内容（让单个决策更鲁棒），而非提高频率去和下层抢时间尺度。
\end{pitfall}

\begin{pitfall}[概念误区：把安全过滤层当成纵向的一环而非横切层]
错误做法：把 CBF 安全过滤放在“MARL $\to$ MPC $\to$ CBF $\to$ WBC”的串联链条里，认为它只过滤 MPC 的输出。后果：MARL 经旁路通道发的命令、或 MPC 内部已违反安全的中间决策，没有被过滤层覆盖，安全保证出现漏洞。根本原因：没理解安全过滤的本质是“控制下发到执行器前的最后一道硬检查”。正确做法：把安全过滤层定位在“所有控制信号汇聚、即将下发到底层执行器”的咽喉点上，保证\textbf{任何}进入底层的控制都经过它（\cref{sec:mh-safety}）。
\end{pitfall}

\begin{pitfall}[编程陷阱：层间用同步阻塞调用而非异步消息]
错误做法：把四层写成一个大循环里的四个同步函数调用（\verb@cmd=marl(); ref=mpc(cmd); safe=cbf(ref); tau=wbc(safe)@），整个循环跑在一个线程里。后果：整个系统的频率被最慢的层（MARL 推理或 MPC 求解）拖到 $10$ Hz，底层 WBC 无法以 $1000$ Hz 运行，执行精度和稳定性崩溃。根本原因：同步串联强制所有层同频，把多频率系统压成了单频率系统。正确做法：各层跑在独立的实时线程/进程，用无锁消息队列通信，每层按自己的频率运行，下层采样上层的“最新可用值”。
\end{pitfall}

\begin{practice}
\begin{enumerate}
  \item \textbf{【架构设计】} 给定“$3$ 台四足机器人协同搬运一根长杆”的任务，画出完整四层架构图，为每层写出 (a) 输入、(b) 输出（具体物理量）、(c) 运行频率、(d) 方法。特别地，思考高层（Level 3）该输出什么命令——是各自的目标位置，还是相对于长杆的抓持点偏移？为什么后者可能更好？
  \item \textbf{【降级设计】} 为练习 1 的系统设计完整的“分层降级”逻辑：分别说明 (a) MARL 推理超时、(b) MPC 求解 infeasible、(c) 某台机器人通信中断 这三种故障下每一层应如何优雅降级，使“长杆不被摔落、机器人不相撞”。
  \item \textbf{【时间尺度分析】} 假设 MPC 求解平均耗时 $25$ ms（最高 $40$ Hz），而你希望底层 WBC 跑 $500$ Hz。计算：(a) 频率比是多少？是否满足时间尺度分离？(b) 两次 MPC 更新之间 WBC 用同一参考运行多少次？(c) 若 MPC 偶尔耗时飙到 $50$ ms 会发生什么？你会如何在 WBC 端平滑这种“卡顿”（参考插值 / 保持上一参考）？
\end{enumerate}
\end{practice}

% ════════════════════════════════════════════════════════════════
\section{RL 高层 + MPC 底层：接口设计}
\label{sec:mh-interface}

\paragraph{动机：一根箭头背后的魔鬼细节。}
\cref{fig:mh-stack} 里 Level 3（MARL）和 Level 2（MPC）之间有一根箭头标着 $u_{\text{cmd}}$。这一节专门讲这根箭头——\textbf{RL 输出的高层命令，到底是什么物理量？MPC 怎么把它变成自己的参考/代价？} MARL 决定“某机向左偏 $0.5$ 米”，MPC 要把它执行出来，听起来简单，但魔鬼藏在细节里——\textbf{“向左偏 $0.5$ 米”这个命令，MPC 到底该怎么用？} 有好几种完全不同的解释：(A) 当成 MPC 的\emph{终端目标位置}（在 $x_H$ 处惩罚离目标点的距离）；(B) 当成\emph{参考轨迹}（生成平滑过渡轨迹逐点跟踪）；(C) 当成代价函数里某项的\emph{权重调整}；(D) 当成一个\emph{约束}（要求轨迹经过某区域）。这四种解释会导致\emph{完全不同的机器人行为}。更要命的是，MARL 是在某一种解释下训练的——如果训练时 MPC 用解释 A、部署时却用解释 B，那么 MARL 学到的策略和 MPC 的实际响应\textbf{对不上}，整个系统行为崩坏。这就是“接口语义对齐”问题。

\begin{insight}[混合架构的成败，$90\%$ 取决于接口设计]
分层混合的真正难点不在“RL 怎么训”或“MPC 怎么解”——这两件事各自都有成熟方案。难点在\textbf{两者之间那个窄接口的语义设计}。接口设计得好，RL 学到的策略能被 MPC 忠实执行，sim-to-real gap 被局部化；设计得差，RL 和 MPC 互相误解，系统在仿真里都未必稳。可以说，\textbf{混合架构的成败，$90\%$ 取决于接口设计，$10\%$ 取决于各层算法本身}。
\end{insight}

\paragraph{反面：三个真实会发生的接口失败模式。}
\begin{itemize}
  \item \textbf{失败模式一：命令空间维度爆炸，RL 学不动。} 初学者常让 MARL 直接输出\emph{整条参考轨迹}（如未来 $2$ 秒、每 $50$ ms 一个点、每点 $6$ 维位姿 $=240$ 维输出），结果动作空间高达几百维，探索效率极低，训练根本收敛不了。正确做法是让 MARL 只输出\emph{低维抽象意图}（如 $3$ 维目标速度），把“意图 $\to$ 轨迹”的展开交给有模型的 MPC。
  \item \textbf{失败模式二：命令语义对动力学敏感，gap 没被局部化。} 让 MARL 输出“期望加速度”或“期望力”，这些量\emph{强依赖动力学}——仿真里产生 $1$ m/s$^2$ 加速度需要的命令和真机不一样（质量、阻力不同），于是 sim-to-real gap 又渗回接口里。正确做法是让命令尽量是\emph{运动学层面的抽象意图}（目标位置/速度/编队偏移），这些量对动力学细节不敏感，“去 A 点”这个意图在仿真和真机里含义完全一样。
  \item \textbf{失败模式三：命令不在 MPC 的“可执行能力”范围内。} 若 MARL 输出一个 MPC 根本做不到的命令（瞬间反向、超出推力极限的加速），MPC 要么 infeasible，要么输出饱和的剧烈控制。正确做法是\emph{约束 MARL 的命令空间到 MPC 的可执行包络内}（如目标速度限幅），或在接口处做软化。
\end{itemize}

\paragraph{历史：从“RL 出动作”到“RL 出命令”的认知演进。}
早期 RL$+$控制混合（2016--2018）常让 RL 直接输出底层动作（力矩/速度），把控制器当成简单限幅或滤波。但人们很快发现这样没有充分利用模型——RL 还是要独自承受大部分 sim-to-real gap。转折点是\emph{分层 RL}（hierarchical RL）思想与控制的结合：HRL 里“高层策略出子目标、低层策略达成子目标”的思想被迁移到“高层 RL 出命令、低层 MPC 达成命令”——区别只是低层从“另一个学习策略”换成“一个模型 MPC”，带来巨大好处（低层 MPC 有保证、有模型、不需训练）。RoCo\cite{mandi2024roco}（2024）把这个思想推到极致——高层甚至不是策略网络而是 LLM；而 MPC $\times$ predictive RL\cite{bjelonic2023mpcrl}（2023）走了另一条路：RL 不出“位置命令”，而是出 MPC 的\emph{尾代价 Q 函数}，让 RL 学“超出 MPC 预测时域之后的长期价值”，补上 MPC“看不远”的短板。这条主线是：\textbf{RL 的输出越来越“高层、抽象、对动力学不敏感”，把越来越多“模型能干的活”还给模型}。

\subsection{理论：命令空间设计的三原则与三种接口形态}
\label{sec:mh-cmd-space}

把上面的失败模式提炼成正面的设计原则。一个好的 RL$\to$MPC 命令接口，其命令空间应满足三条（\cref{tab:mh-cmd-principle}）。

\begin{table}[H]\centering\small
\caption{命令空间设计三原则}
\label{tab:mh-cmd-principle}
\begin{tabular}{p{2.4cm}p{5.6cm}p{4.2cm}}
\toprule
原则 & 内容 & 违反的后果 \\
\midrule
低维性 & 命令维度尽量低（通常 $\le 6$ 维/agent），把“展开成轨迹”留给 MPC & 维度爆炸，RL 训练不收敛 \\
运动学抽象性 & 命令是运动学意图（目标位置/速度/编队偏移），不是动力学量（力/加速度） & sim-to-real gap 渗回接口，迁移失败 \\
可执行性 & 命令空间约束在 MPC 的可执行包络内（限幅 / 投影到可行集） & MPC 频繁 infeasible 或剧烈饱和 \\
\bottomrule
\end{tabular}
\end{table}

基于三原则，三种典型命令接口形态：
\begin{description}
  \item[形态一：目标速度接口（velocity command）] 最常用，适合移动/编队。MARL 输出每个 agent 的期望质心速度 $v_{\text{cmd},i}\in\mathbb{R}^{2}$ 或 $\mathbb{R}^{3}$；MPC 把它当参考速度，代价函数惩罚 $\|v-v_{\text{cmd}}\|^2$，在动力学和约束下跟踪。优点：维度极低（$2$--$3$ 维）、运动学抽象、对动力学不敏感、天然限幅（限制在最大速度内）。
  \item[形态二：编队偏移接口（formation offset）] 适合保持编队的协同。MARL 输出每个 agent 相对编队中心（或虚拟领队）的偏移 $\Delta p_i$；MPC 让编队中心按预设轨迹走，每个 agent 的目标位置 $=$ 中心位置 $+\Delta p_i$。优点：把“队形”和“整体运动”解耦，MARL 只学“队形怎么动态调整”（避障时压扁、过通道时排成一列）。
  \item[形态三：子目标 / 航点接口（subgoal / waypoint）] 适合长程导航 $+$ 局部规划。MARL 输出每个 agent 的下一个子目标位置 $g_i$（一个相对较远的航点）；MPC（或局部规划器）负责从当前位置安全到达 $g_i$、处理局部避障。优点：MARL 做长程、组合的“去哪”决策，MPC 做局部、连续的“怎么到”，分工最清晰。
\end{description}

\begin{note}
\textbf{类比（标边界）：目标速度接口 $\approx$ 给车装了自适应巡航（ACC）.}\;
司机（MARL）只用油门表达“想多快”，由 ACC（MPC）去管油门开度、换挡、跟车距离。\emph{相似}：司机只下达高层意图，底层系统负责把意图翻译成具体、可行、安全的执行细节。\emph{不同}：ACC 的“翻译”是固定规则，而 MPC 的翻译是带完整动力学模型的在线优化；且 ACC 只管纵向速度，目标速度接口是多维（含方向）的协调命令。\textbf{类比的边界}：不要延伸到“司机和 ACC 各自独立”——在我们的系统里 MARL 是被训练得专门适配某个特定 MPC 的（接口语义训练部署须一致），而司机换辆车照样能开。
\end{note}

\subsection{理论-工程桥接：最小可运行的 RL$\to$MPC 接口（目标速度形态）}
\label{sec:mh-iface-code}

把“目标速度接口”落地。核心是把策略输出的归一化动作 $a\in[-1,1]^3$ 翻译成 MPC 能直接消费的参考速度 $v_{\text{cmd}}$，且保证它落在 MPC 可执行包络内：(1) 归一化到物理速度（运动学抽象量，对动力学不敏感）；(2) 限幅到 $v_{\max}$（可执行性原则，绝不发 MPC 做不到的命令）；(3) 限制相邻周期的速度变化率 $\le a_{\max}\Delta t$（避免命令跳变导致 MPC 输出剧烈饱和）。下面的最小实现用 \verb@policy_action_to_vcmd@ 完成这三步，再用 \verb@build_mpc_cost@ 把 $v_{\text{cmd}}$ 放进 MPC 代价的速度跟踪项。

\begin{codebox}[Python]{RL$\to$MPC 目标速度接口（归一化动作 $\to$ 安全可执行参考速度）}
import numpy as np

class RLToMPCInterface:
    def __init__(self, v_max, a_max, ctrl_dt):
        self.v_max, self.a_max, self.ctrl_dt = v_max, a_max, ctrl_dt
        self.v_cmd_prev = np.zeros(3)              # last cmd, for smoothing

    def policy_action_to_vcmd(self, action):       # action in [-1,1]^3
        v_target = action * self.v_max             # 1) normalized -> physical speed
        speed = np.linalg.norm(v_target)           # 2) clamp to executable envelope
        if speed > self.v_max:
            v_target = v_target * (self.v_max / speed)
        dv = v_target - self.v_cmd_prev            # 3) rate limit (protect MPC)
        dv_max = self.a_max * self.ctrl_dt
        if np.linalg.norm(dv) > dv_max:
            dv = dv * (dv_max / np.linalg.norm(dv))
        self.v_cmd_prev = self.v_cmd_prev + dv
        return self.v_cmd_prev                      # 3-D, fed to MPC cost

# v_cmd enters the MPC COST (soft), never a hard constraint:
#   min  sum_k ||v_k - v_cmd||^2_{Qv} + ||u_k||^2_R
#   s.t. x_{k+1}=f(x_k,u_k),  u_k in U,  g(x_k,u_k)<=0   (hard: friction cone/limits)
# so even if v_cmd is momentarily unreachable, MPC stays feasible (best-effort track).
\end{codebox}

\noindent\textbf{三处易错（务必避开）}：(i) RL 直接输出整条参考轨迹 $\to$ 动作空间几百维、训练不收敛；应只输出 $3$ 维意图、让 MPC 展开。(ii) RL 输出动力学量（期望力）$\to$ gap 渗回接口；应输出运动学意图。(iii) 把 $v_{\text{cmd}}$ 当 MPC \emph{硬等式约束}（\verb@v_k == v_cmd@）$\to$ 命令暂时不可达时 MPC 直接 infeasible 无解；命令应进\textbf{代价}（软），硬约束只留给真正不可违反的物理约束。

\subsection{一个进阶接口：RL 出 MPC 的尾代价}
\label{sec:mh-terminal-cost}

除了“RL 出命令、MPC 跟踪”这种主流形态，还有一种更精巧的接口——\textbf{RL 学习 MPC 的尾代价}（MPC $\times$ predictive RL\cite{bjelonic2023mpcrl}）。MPC 的一个根本局限是\emph{预测时域有限}（只看未来 $H$ 步），它对“$H$ 步之外会发生什么”一无所知，只能用一个手工设计的尾代价 $V_f(x_H)$ 近似“末端状态的长期价值”。但这个长期价值恰恰\emph{难以手工设计}——它本质上是一个最优值函数，而值函数正是 RL 的看家本领。于是这个接口让 \textbf{RL 学一个 Q 函数当 MPC 的尾代价}：
\begin{equation}
  \min_{u_0,\dots,u_{H-1}}\ \sum_{k=0}^{H-1}\ell(x_k,u_k)+\underbrace{Q_\theta(x_H,u_H)}_{\text{RL 学的尾代价}},\quad\text{s.t. 动力学 $+$ 约束.}
  \label{eq:mh-terminal-q}
\end{equation}
这样 MPC 的短时域优化“接上了”RL 学到的长期价值，既保留 MPC 的硬约束保证和短期精确性，又通过 RL 的尾代价获得长程全局视野。这是“RL 补 MPC 看不远的短板”的精妙体现——和“RL 出命令”形成互补：前者 RL 在\emph{时间维度}上补 MPC（看更远），后者 RL 在\emph{决策维度}上补 MPC（出难建模的命令）。

\begin{insight}[RL$\to$MPC 的接口可以建在不同的“语义位置”上]
RL$\to$MPC 的接口可以是 MPC 的\textbf{输入}（命令/参考），也可以是 MPC 的\textbf{代价}（尾代价 Q 函数，\cref{eq:mh-terminal-q}），甚至是 MPC 的\textbf{约束}或\textbf{模型残差}（\cref{sec:mh-residual}）。每个位置对应一种“让 RL 补 MPC 不同短板”的方式：补“难建模的决策”（命令）、补“看不远的长期价值”（尾代价）、补“模型误差”（残差）。理解这一点，你就不会把“RL$+$MPC 混合”当成固定模板，而会根据“MPC 到底缺什么”来选择接口的接入点。
\end{insight}

\begin{pitfall}[思维陷阱：让 RL 直接输出底层动作，浪费 MPC 的模型能力]
错误做法：让 MARL 直接输出关节力矩或整条参考轨迹，MPC/控制器只做简单限幅。后果：RL 动作空间维度过高训练不收敛，或 RL 独自承受全部 sim-to-real gap 迁移失败，MPC 的模型优势完全没用上。根本原因：RL 出底层动作等于把 MPC 降级成一个滤波器，丢掉它最值钱的能力——基于模型的约束满足和轨迹生成。正确做法：RL 只输出低维运动学意图（$\le 6$ 维），把“意图$\to$可行控制”完全交给 MPC。
\end{pitfall}

\begin{pitfall}[概念误区：训练和部署用不同的命令语义，导致 RL 和 MPC 对不上]
错误做法：训练时 MPC 把 $v_{\text{cmd}}$ 当终端目标位置（解释 A），部署时却当参考速度跟踪（解释 B）。后果：MARL 学到的策略和 MPC 的实际响应行为对不上，仿真训练好的策略部署后行为诡异。根本原因：MARL 是在某种“命令$\to$行为”映射下学的，这个映射由 MPC 如何解释命令决定，换了解释 $=$ 换了环境。正确做法：在训练\textbf{之前}就冻结接口语义（命令是什么物理量、MPC 怎么解释、限幅平滑规则），训练和部署用\emph{完全相同}的接口代码，把接口当成环境的一部分。
\end{pitfall}

\begin{pitfall}[编程陷阱：$v_{\text{cmd}}$ 进 MPC 硬约束导致 infeasible]
错误做法：在 MPC 里写 \texttt{v\_k == v\_cmd} 作为硬等式约束。后果：命令因障碍或动力学暂时不可达时，MPC 求解 infeasible，控制输出中断，系统失控。根本原因：混淆了“命令是期望”和“命令是约束”。正确做法：$v_{\text{cmd}}$ 永远进\textbf{代价函数}（软跟踪 $\|v-v_{\text{cmd}}\|^2_{Q_v}$），保证 MPC 在任何命令下都有解；硬约束只留给真正不可违反的物理约束。
\end{pitfall}

\begin{practice}
\begin{enumerate}
  \item \textbf{【接口实现】} 基于 \verb@RLToMPCInterface@，实现一个“编队偏移接口”：MARL 为每个 agent 输出归一化偏移动作，接口翻译成相对编队中心的偏移 $\Delta p_i$，并保证 (a) 偏移幅值不超过 \verb@max_offset@（避免编队散架），(b) 任意两 agent 的目标位置间距不小于 \verb@min_sep@（避免目标本身冲突）。说明 (b) 这个约束为什么应在接口层做、而不是留给下游 CBF。
  \item \textbf{【接口调试】} 给定现象“仿真里训练好的 $8$ 机编队策略，部署到真机后整体能协调，但每台机器人都在高频抖动”。列出\textbf{至少三个}可能的接口层根因（从本节失败模式和陷阱里找），并为每个给出排查方法和修复方案。
  \item \textbf{【设计扩展】} 设计第三种接口——\textbf{RL 动态调整 MPC 代价函数的权重}（如 RL 输出“此刻应更看重避障还是更看重跟踪”的权重比）。写出：(a) RL 的输出是什么、维度多少；(b) MPC 端如何消费这个权重；(c) 相比“RL 出命令”的优劣（提示：它对哪类任务有独特价值，以及会不会让 MPC 行为变得难以预测）。
\end{enumerate}
\end{practice}

% ════════════════════════════════════════════════════════════════
\section{安全过滤层：CBF-QP 与 GCBF+}
\label{sec:mh-safety}

\paragraph{动机：给黑盒策略套一个“安全外壳”。}
\cref{sec:mh-why} 反复强调：MARL 的奖励惩罚是软的，无法保证安全。那怎么办？一个绝妙的思路是——\textbf{不要试图让策略本身安全，而是在策略和执行器之间插一道“安全过滤器”}。上层（MARL 或 MPC）随便输出一个名义控制 $u_{\text{nom}}$，安全过滤层在它执行前做一次\emph{最小修正}，把它投影到“保证不违反硬约束”的安全集里，得到 $u_{\text{safe}}$。这个思路的革命性在于：\textbf{它把“安全”和“性能”彻底解耦了}。MARL 全力优化性能（协调、效率），完全不用操心安全；安全过滤层全力保证安全，完全不用操心性能。两个目标在两个独立模块里、用两种不同的数学（RL 的期望优化 vs CBF 的逐状态约束）各自处理——这正是 \cref{sec:mh-why} “安全必须是约束而非目标项”那个洞察的工程落地。

\begin{note}
\textbf{类比（标边界）：安全过滤 $\approx$ 自动紧急制动（AEB）.}\;
新手司机（MARL 策略）想怎么开就怎么开，但只要 AEB（安全过滤器）检测到“再这样下去要撞了”，它就\emph{最小限度地介入}——能不踩刹车就不踩，一旦要撞就强制干预，且只干预到“刚好不撞”为止。\emph{相似}：两者都是“让主控制器自由发挥，只在危险时最小介入”。\emph{不同}：AEB 是基于规则的离散触发（撞前一刻全力刹车），而 CBF 过滤是\textbf{连续、平滑、逐时刻}的修正（持续地、温和地把控制往安全方向掰），且有\emph{数学可证明}的安全保证。\textbf{类比的边界}：不要把 CBF 想成“危险时才触发”，它是\textbf{时刻都在}工作的。
\end{note}

\paragraph{反面：不做安全过滤只剩两个糟糕选择。}
\textbf{选择一：信任 MARL 的软安全。} 这等于接受一个非零的碰撞概率，真机上迟早出事；且这个概率在 OOD 情况下会突然飙升（训练时 $0.1\%$，遇到没见过的风扰可能瞬间变成 $50\%$），你完全无法预测什么时候出事。\textbf{选择二：把所有安全约束塞进 MPC 硬约束。} 这看起来可行，但 (1) MPC 是\emph{有限时域}的，只保证未来 $H$ 步安全，$H$ 步之外不管；(2) 多机器人的 pairwise 碰撞约束数量随 $N^2$ 增长，全塞进一个 MPC 会让 QP 规模爆炸、求解超时；(3) 分布式时每个 agent 的 MPC 看不到全局，难以协调避碰。CBF 过滤恰好补上这三点：它保证\textbf{无限时域}的前向不变性（不是只看 $H$ 步），它\emph{计算极轻}（每个 agent 一个小 QP），它\emph{天然分布式}（每个 agent 只用局部邻居信息）。

\paragraph{历史：从 CBF 到 GCBF+。}
控制屏障函数（CBF）的现代形式由 Ames 等人在 2014--2019 年间确立\cite{ames2017cbf,ames2019cbf}。它的思想源头是控制李雅普诺夫函数（CLF）——CLF 用一个能量函数保证\emph{稳定性}，CBF 用一个屏障函数保证\emph{安全性}。Ames 的关键贡献是把“安全”写成一个\emph{仿射于控制 $u$ 的不等式约束}，从而可以用一个二次规划（QP）在每个控制周期实时求解——这就是 CBF-QP，一个可叠加在任何名义控制器之上的“安全过滤器”。把它推广到\emph{多机器人碰撞避免}的奠基工作是 Wang--Ames--Egerstedt\cite{wang2017safety}；而\cref{ch:plan-safemarl} 已论证：当每个机器人独立解自己的最小侵入 CBF-QP 时，它们集体达到的状态恰好是一个广义 Nash 均衡——“安全证书”自动成为“博弈均衡”的一部分。本章承接其结论、聚焦“把它作为混合系统的安全过滤层”，重叠处一律指过去（\cref{ch:plan-safemarl}）。但经典 CBF 有一个致命的可扩展性问题：\textbf{手工设计 CBF 函数 $h$ 对多智能体几乎不可行}。GCBF+\cite{zhang2025gcbfplus}（Zhang, Fan 等，T-RO 2024）的突破正在于此：\textbf{用图神经网络（GNN）参数化一个 CBF}，让它天然可扩展到任意数量的 agent。

\subsection{理论：CBF 与前向不变性（先复习）}
\label{sec:mh-cbf-recall}

我们先把 CBF 的核心快速复活一遍（详证见 \cref{ch:plan-safemarl}）。考虑一个\textbf{控制仿射系统} $\dot x=f(x)+g(x)u$，其中 $f(x)$ 是漂移项、$g(x)$ 是输入项。定义\textbf{安全集}为某个函数 $h(x)$ 的超水平集：
\begin{equation}
  \mathcal{C}=\{x:h(x)\ge0\},\quad \partial\mathcal{C}=\{x:h(x)=0\},\quad \mathrm{Int}(\mathcal{C})=\{x:h(x)>0\}.
  \label{eq:mh-safe-set}
\end{equation}
直觉上 $h(x)$ 是一个“到危险的距离”——$h>0$ 安全、$h=0$ 危险边缘、$h<0$ 已违反。比如碰撞避免里 $h=\|p_i-p_j\|^2-(2r)^2$（两 agent 距离平方减安全半径平方）。

\begin{definition}[控制屏障函数与 CBF 条件]\label{def:mh-cbf}
$h$ 是一个控制屏障函数（CBF），若存在一个扩展 class-$\mathcal{K}$ 函数 $\alpha$（严格增、$\alpha(0)=0$），使得对所有 $x\in\mathcal{C}$ 存在控制 $u$ 满足
\begin{equation}
  \dot h(x,u)=\underbrace{L_fh(x)}_{\frac{\partial h}{\partial x}f(x)}+\underbrace{L_gh(x)}_{\frac{\partial h}{\partial x}g(x)}\,u\ \ge\ -\alpha(h(x)),
  \label{eq:mh-cbf-cond}
\end{equation}
其中 $L_fh,L_gh$ 是 $h$ 沿漂移项/输入项的李导数\cite{ames2017cbf}。
\end{definition}

\begin{theorem}[前向不变性]\label{thm:mh-fwd-invariance}
若 $x(0)\in\mathcal{C}$ 且每时刻控制满足 CBF 条件 \cref{eq:mh-cbf-cond}，则 $x(t)\in\mathcal{C}$ 对所有 $t\ge0$ 成立——即安全集 $\mathcal{C}$ 前向不变。
\end{theorem}
\begin{proof}[证明思路]\rebuilt{（据 \cite{ames2017cbf}，完整证见 \cref{ch:plan-safemarl}）}
关键在边界附近的行为。当状态逼近安全边界时 $h\to0^+$，于是 $\alpha(h)\to0$，\cref{eq:mh-cbf-cond} 变成 $\dot h\ge0$——它\emph{强制 $h$ 不再下降}，把状态“挡”在边界内侧；而在内部深处（$h$ 大）$\alpha(h)$ 大、约束很宽松（允许 $\dot h$ 取较大负值即朝边界移动），但不会越界。形式上，比较引理给出 $h(x(t))\ge\beta(h(x(0)),t)\ge0$，其中 $\beta$ 是由 $\dot h=-\alpha(h)$ 定义的 class-$\mathcal{KL}$ 函数。
\end{proof}

\begin{insight}[$\alpha(h)$ 是一个“自适应的放松率”]
CBF 不等式 $\dot h\ge-\alpha(h)$ 的精妙之处在于 $\alpha(h)$ 这个自适应的放松率：它让约束在安全深处\textbf{宽松}（不必要地限制控制是浪费）、在边界附近\textbf{收紧}（必须阻止越界）。如果把 $-\alpha(h)$ 换成固定常数（如 $\dot h\ge0$，要求 $h$ 永不下降），系统会过度保守——明明离危险很远也不许靠近。$\alpha(h)$ 用一个随距离自适应的边界，在“安全”和“不保守”之间取得了最优平衡——这是 CBF 远胜于简单“距离阈值刹车”的根本原因。
\end{insight}

\paragraph{相对度与高阶 CBF（HOCBF）。}
\cref{def:mh-cbf} 的推导假设 $L_gh\ne0$，即 $\dot h$ \emph{显含} $u$。但很多情况下不是这样——比如 $h$ 是位置的函数（$h=\|p_i-p_j\|^2-(2r)^2$），而控制 $u$ 是加速度/力，那么 $\dot h$ 只含速度、不含 $u$，要再求一次导 $\ddot h$ 才出现 $u$。这种“$h$ 对 $u$ 的\emph{相对度}（relative degree）$\ge2$”的情况需要\textbf{高阶 CBF}（high-order CBF, HOCBF\cite{xiao2022hocbf}）：构造一串嵌套函数 $\psi_0=h$，$\psi_1=\dot\psi_0+\alpha_1(\psi_0)$，$\psi_2=\dot\psi_1+\alpha_2(\psi_1)$，直到 $\psi_k$ 显含 $u$，然后约束 $\dot\psi_{k-1}\ge-\alpha_k(\psi_{k-1})$。这是多机器人（二阶/三阶动力学）必然遇到的——力控机器人的碰撞避免几乎都是相对度 $2$ 的问题，必须用 HOCBF。

\subsection{理论：CBF-QP 安全过滤器的推导}
\label{sec:mh-cbf-qp}

现在把 CBF 变成一个“过滤器”。给定一个名义控制 $u_{\text{nom}}$（来自 MARL 或 MPC），我们要找一个\emph{尽量接近 $u_{\text{nom}}$、但满足 CBF 安全约束}的 $u_{\text{safe}}$。这自然写成一个二次规划：
\begin{equation}
  u_{\text{safe}}=\arg\min_{u}\ \tfrac12\|u-u_{\text{nom}}\|^2\quad\text{s.t.}\quad L_fh(x)+L_gh(x)\,u\ge-\alpha(h(x)).
  \label{eq:mh-cbf-qp}
\end{equation}
目标函数 $\|u-u_{\text{nom}}\|^2$ 是“最小修正”——尽量少改动名义控制（\emph{当名义控制本就安全时，不做任何修改}）；约束是 CBF 安全条件。这是一个\textbf{单约束的 QP}，对单个安全约束有\emph{闭式解}，我们逐步推导它。

把约束改写成标准形式 $a^\top u\le b$，其中 $a=-L_gh(x)^\top$、$b=L_fh(x)+\alpha(h(x))$（移项：$L_gh\cdot u\ge-\alpha(h)-L_fh$ 即 $-L_gh\cdot u\le L_fh+\alpha(h)$）。问题变成 $\min_u\tfrac12\|u-u_{\text{nom}}\|^2$ s.t.\ $a^\top u\le b$。

\begin{theorem}[CBF-QP 单约束闭式解]\label{thm:mh-cbf-closed}
\cref{eq:mh-cbf-qp} 的解为
\begin{equation}
  u_{\text{safe}}=u_{\text{nom}}-\frac{\max\!\big(0,\ a^\top u_{\text{nom}}-b\big)}{\|a\|^2}\,a,\qquad a=-L_gh^\top,\ \ b=L_fh+\alpha(h).
  \label{eq:mh-cbf-closed}
\end{equation}
\end{theorem}
\begin{proof}
\textbf{情况一：$u_{\text{nom}}$ 已满足约束}（$a^\top u_{\text{nom}}\le b$）。约束不激活，最优解就是 $u_{\text{safe}}=u_{\text{nom}}$——安全控制时过滤器透明、零修正。\textbf{情况二：$u_{\text{nom}}$ 违反约束}（$a^\top u_{\text{nom}}>b$）。约束激活（取等号），用 Lagrangian 求解。构造 $\mathcal{L}=\tfrac12\|u-u_{\text{nom}}\|^2+\lambda(a^\top u-b)$，$\lambda\ge0$。KKT 稳定性条件：
\begin{equation}
  \frac{\partial\mathcal{L}}{\partial u}=(u-u_{\text{nom}})+\lambda a=0\ \Rightarrow\ u=u_{\text{nom}}-\lambda a.
  \label{eq:mh-kkt-stat}
\end{equation}
代入激活约束 $a^\top u=b$：$a^\top(u_{\text{nom}}-\lambda a)=b\Rightarrow\lambda=\dfrac{a^\top u_{\text{nom}}-b}{\|a\|^2}$。两种情况合并即得 \cref{eq:mh-cbf-closed}（$\max(0,\cdot)$ 保证只在违反时修正）。
\end{proof}

\begin{insight}[安全过滤的几何本质是“到安全半空间的正交投影”]
闭式解 \cref{eq:mh-cbf-closed} 揭示安全过滤的几何本质——它就是一个\textbf{到安全半空间的欧氏投影}。安全集在控制空间里是一个半空间 $\{u:a^\top u\le b\}$；过滤器做的事，就是把落在外面的名义控制“拉”回到最近的边界点（$\max(0,\cdot)$ 保证只在违反时修正，$\frac{\cdot}{\|a\|^2}a$ 是投影到超平面的最小位移）。这解释了为什么安全过滤几乎不损害性能：投影是“最小改动”，它只在你真要越界时、且只把你拉回到刚好不越界的位置。性能损失被压到了数学上的最小值。
\end{insight}

\paragraph{一个手算的数值例子（让闭式解从抽象变具体）。}
考虑最简单的一阶 integrator 单机避障：$\dot x=u$（$x,u\in\mathbb{R}^2$），机器人在原点 $x=(0,0)$，要避开位于 $p_{\text{obs}}=(2,0)$、半径 $r=1$ 的圆形障碍。安全函数 $h(x)=\|x-p_{\text{obs}}\|^2-r^2$，当前 $h=\|(0,0)-(2,0)\|^2-1=4-1=3>0$（安全）。取 $\alpha(h)=h$（线性 $\mathcal{K}$ 类函数）。名义控制 $u_{\text{nom}}=(3,0)$——MARL 想让机器人\emph{直接冲向障碍}（正 $x$ 方向，速度 $3$）。
\begin{itemize}
  \item 对一阶 integrator，$f(x)=0$，$g(x)=\mathbf{I}$，所以 $L_fh=0$，$L_gh=\frac{\partial h}{\partial x}=2(x-p_{\text{obs}})=2(0-2,0-0)=(-4,0)$。
  \item 约束系数 $a=-L_gh^\top=(4,0)^\top$，$b=L_fh+\alpha(h)=0+3=3$。
  \item 检查名义控制是否违反：$a^\top u_{\text{nom}}=(4,0)\cdot(3,0)=12$，而 $b=3$。$12>3$ $\to$ \emph{违反约束}，需修正。
  \item 代入闭式解：$\lambda=\dfrac{a^\top u_{\text{nom}}-b}{\|a\|^2}=\dfrac{12-3}{16}=\dfrac{9}{16}=0.5625$。
  \item $u_{\text{safe}}=u_{\text{nom}}-\lambda a=(3,0)-0.5625\cdot(4,0)=(3-2.25,\ 0)=(0.75,\ 0)$。
\end{itemize}
\textbf{解读}：过滤器没有把机器人完全拦停（$u_{\text{safe}}$ 仍是 $(0.75,0)$，还在朝障碍方向走），而是把朝障碍的速度从 $3$ \emph{减到 $0.75$}——刚好让 $\dot h=-\alpha(h)$ 这个边界条件满足（验证：$\dot h=L_gh\cdot u_{\text{safe}}=(-4,0)\cdot(0.75,0)=-3=-\alpha(h)=-h$ $\checkmark$）。这正是“最小修正”的精髓：它不过度干预（没有强制停止或反向），只把违反的分量削减到刚好不越界。随着机器人继续靠近、$h$ 减小，$b=\alpha(h)$ 也减小，允许的朝向速度持续下降，最终在边界处降到零——平滑、连续、绝不越界。

\paragraph{多个安全约束的情况。}
多机器人系统里一个 agent 通常对多个邻居都有碰撞约束，于是有多个 CBF 不等式同时约束 $u$：
\begin{equation}
  \min_u\tfrac12\|u-u_{\text{nom}}\|^2\quad\text{s.t.}\quad Au\le b,\quad A=[a_1^\top;a_2^\top;\dots],\ b=[b_1;b_2;\dots].
  \label{eq:mh-cbf-multi}
\end{equation}
多约束 QP 一般没有简单闭式解（除非只有一个约束激活），需用数值 QP 求解器（如 OSQP、qpOASES）。但它仍极轻量——典型 agent 的邻居数有限（几个到十几个），QP 规模很小，单个 agent 在亚毫秒级就能求解。

\subsection{理论-工程桥接：用 OSQP 实现 pairwise CBF 碰撞避免（HOCBF）}
\label{sec:mh-cbf-code}

把理论落地。场景：$N$ 个 2D 双积分器（double integrator）agent，控制是加速度，要保证任意两 agent 永不碰撞。注意双积分器里 $h$（位置函数）对 $u$（加速度）的相对度是 $2$，所以\textbf{必须用 HOCBF}：因为 $\ddot x=u$，$h=\|p_i-p_j\|^2-(2r)^2$，$\dot h=2(p_i-p_j)\cdot(v_i-v_j)$ 只含速度、不含 $u$；构造 $\psi_1=\dot h+\alpha_1h$，再约束 $\dot\psi_1\ge-\alpha_2\psi_1$，展开后 $\dot\psi_1$ 含 $\ddot h=2\|v_{ij}\|^2+2p_{ij}\cdot(u_i-u_j)$ 即含 $u$，这才“逼出”了对 $u$ 的约束。下面的实现把每对 $(i,j)$ 的 HOCBF 约束写成 $Au\le b$ 喂给 OSQP。

\begin{codebox}[Python]{pairwise HOCBF 碰撞避免安全过滤（OSQP 后端）}
import numpy as np, osqp
from scipy import sparse

def pairwise_cbf_filter(p, v, u_nom, r_safe, alpha1=2.0, alpha2=2.0):
    # N double integrators (acc control). h=||pi-pj||^2-(2r)^2, rel-deg 2 -> HOCBF.
    N = p.shape[0]; nu = 2 * N
    P = sparse.eye(nu, format='csc'); q = -u_nom.flatten()   # min 1/2||u-u_nom||^2
    A_rows, b_vals = [], []
    for i in range(N):
        for j in range(i + 1, N):
            pij, vij = p[i] - p[j], v[i] - v[j]
            h = pij @ pij - (2 * r_safe) ** 2          # h
            h_dot = 2.0 * pij @ vij                     # hdot (no u)
            psi1 = h_dot + alpha1 * h                   # psi1
            # HOCBF: psidot1 >= -a2 psi1, with hddot = 2||vij||^2 + 2 pij.(ui-uj)
            #   -2 pij.ui + 2 pij.uj <= 2||vij||^2 + a1 hdot + a2 psi1
            a_row = np.zeros(nu)
            a_row[2*i:2*i+2] = -2.0 * pij; a_row[2*j:2*j+2] = +2.0 * pij
            A_rows.append(a_row)
            b_vals.append(2.0 * (vij @ vij) + alpha1 * h_dot + alpha2 * psi1)
    if not A_rows:
        return u_nom
    prob = osqp.OSQP()
    prob.setup(P=P, q=q, A=sparse.csc_matrix(np.array(A_rows)),
               l=-np.inf*np.ones(len(b_vals)), u=np.array(b_vals), verbose=False)
    res = prob.solve()
    if res.info.status_val not in (1, 2):     # infeasible -> SAFE fallback
        return np.zeros_like(u_nom)           # zero accel (conservative but safe)
    return res.x.reshape(N, 2)
\end{codebox}

\noindent\textbf{三处易错}：(i) 对双积分器用一阶 CBF（相对度搞错）$\to$ $\dot h$ 不含 $u$，约束形同虚设、碰撞照发；必须 HOCBF。(ii) QP infeasible 时退回 $u_{\text{nom}}$ $\to$ infeasible 往往恰发生在“危险逼近”时（约束最紧），退回名义控制等于在最危险时刻关闭保护，必撞；应降级到保守安全动作（零加速度/最大减速）。(iii) $\alpha$ 取太大 $\to$ 约束 $\dot h\ge-\alpha h$ 的右端 $-\alpha h$ 是巨大负数，约束几乎总满足、形同虚设，现象是 agent 冲到极近才“突然”刹车、常来不及。

\begin{table}[H]\centering\small
\caption{三种“安全”机制对比：奖励惩罚 vs MPC 硬约束 vs CBF 过滤}
\label{tab:mh-safety-compare}
\begin{tabular}{p{2.6cm}p{2.8cm}p{2.4cm}p{2.2cm}p{2.4cm}}
\toprule
机制 & 保证强度 & 时域 & 计算量 & 可扩展性($N$) \\
\midrule
MARL 奖励惩罚 & 软（期望意义） & --- & 训练期 & 好但无保证 \\
MPC 硬约束 & 硬（但有限时域） & 未来 $H$ 步 & 大（全局 QP） & 差（$N^2$ 约束） \\
CBF-QP 过滤 & 硬（无限时域，前向不变） & 逐时刻 & 极小（局部 QP） & 极好（分布式） \\
\bottomrule
\end{tabular}
\end{table}

\subsection{GCBF+：用 GNN 把安全过滤扩展到上千 agent}
\label{sec:mh-gcbfplus}

pairwise CBF 过滤有一个隐患：\textbf{手工 CBF 在密集场景下会过度保守甚至 infeasible}。当一个 agent 被很多邻居包围时，各 pairwise 约束可能互相冲突（往左躲 A 就撞 B，往右躲 B 就撞 A），QP 无解；手工设计一个能协调所有邻居、永远可行的 CBF，对大 $N$ 几乎不可能。GCBF+\cite{zhang2025gcbfplus} 的解法是\textbf{学一个 CBF}：用 GNN 把“agent 的安全值”参数化为其局部邻居图的函数
\begin{equation}
  h_\theta=\mathrm{GNN}_\theta(\text{node features},\ \text{edge features},\ \text{graph topology}).
  \label{eq:mh-gcbf}
\end{equation}
为什么 GNN 是天作之合？因为多智能体安全问题有几个结构性质，恰好是 GNN 的强项（\cref{tab:mh-gnn-match}）。

\begin{table}[H]\centering\small
\caption{多智能体安全问题的结构性质与 GNN 的匹配}
\label{tab:mh-gnn-match}
\begin{tabular}{p{2.6cm}p{5.0cm}p{4.2cm}}
\toprule
性质 & 含义 & GNN 为何匹配 \\
\midrule
置换等变性 & agent 的安全不依赖编号顺序 & 消息传递天然置换等变 \\
邻居数不变性 & 同一 CBF 要能处理“$3$ 个邻居”和“$30$ 个邻居” & 聚合（aggregation）对邻居数不敏感 \\
局部性 & agent 的安全只依赖近邻 & 有限跳消息传递 $=$ 局部感受野 \\
\bottomrule
\end{tabular}
\end{table}

\noindent GCBF+ 的训练损失同时优化“CBF 的有效性”（违反时惩罚）和“CBF 条件的满足”，三部分各编码 CBF 的一个要求\cite{zhang2025gcbfplus}：
\begin{equation}
  \mathcal{L}_{\text{GCBF+}}=\underbrace{\textstyle\sum\max(0,-h_\theta+\epsilon)\big|_{\text{safe}}}_{\text{安全态应有 }h>0}+\underbrace{\textstyle\sum\max(0,h_\theta+\epsilon)\big|_{\text{unsafe}}}_{\text{危险态应有 }h<0}+\underbrace{\textstyle\sum\max(0,-\dot h_\theta-\alpha(h_\theta))}_{\text{CBF 条件违反惩罚}}.
  \label{eq:mh-gcbf-loss}
\end{equation}
部署时，每个 agent 用学到的 $h_\theta$ 跑一个本地 CBF-QP（和上面的 pairwise 过滤同构，只是 $h$ 从手工公式换成 GNN 输出）。GCBF+ 关键工程亮点：\emph{可扩展}（在 $256$ agent 上比手工 CBF 好 $20\%$，在 $1024$ agent 上比纯 RL 好 $40\%$）；\emph{接受原始感知}（直接吃 LiDAR 点云，不需理想的邻居状态估计——这对真机至关重要）；\emph{考虑执行约束}（用 look-ahead 机制把执行器限制纳入 CBF 训练）。

\begin{note}
\textbf{关键实验数据（数量级，具体以论文为准）.}\;\rebuilt{以下数字转述自 \cite{zhang2025gcbfplus}，待联网核对：中小规模（$\le256$ agents）安全率最高、手工 CBF-QP 约低 $20\%$；大规模（$1024$ agents）比领先 RL（如 MAPPO）约高 $40\%$；训练只用 $8$ 个 agent。\pz{存疑：$256$ agents 高 $20\%$、$1024$ agents 高 $40\%$ 等具体数值需对照论文核实；GCBF（前身）vs GCBF+ 勿张冠李戴。}}
真正卖点是“不牺牲任务的安全”——手工 CBF 把 $\alpha$ 调小也很安全但过保守（绕大圈、到不了），RL 难两全；GCBF+ 用“以 QP 控制器作 reference $+$ 三部分损失”打破“安全 vs 任务”权衡。看安全方法时永远要同时看“安全率”和“任务完成率”。
\end{note}

\begin{insight}[GCBF+ 学的不是策略，而是约束]
GCBF+ 把 CBF 从“手工设计的安全证书”变成“学习得到的安全证书”，但\textbf{它学的不是策略，而是约束}——这是它和纯 RL 的根本区别。纯 RL 学一个“倾向于安全的策略”（软），GCBF+ 学一个“定义安全集边界的函数”（硬，因为部署时仍用 QP 强制满足 CBF 条件）。所以 GCBF+ 既有学习的可扩展性（GNN 处理任意 $N$），又有 CBF 的硬保证（QP 投影）。当然，由于 $h_\theta$ 是学出来的，它的安全保证是\textbf{概率性}的（依赖训练覆盖），如何走向形式化认证仍是开放问题——这是它诚实的局限。
\end{insight}

\paragraph{一个关键的工程权衡：过滤的“保守性”。}
安全过滤不是没有代价的——它有一个根本的权衡轴：\textbf{保守性 vs 任务性能}。$\alpha$ 函数和安全半径 $r$ 的选择直接控制这个权衡：$\alpha$ 小 / $r$ 大 $\to$ 过滤器\emph{保守}（很早就减速避让，绝对安全但任务效率低，密集场景可能完全卡死）；$\alpha$ 大 / $r$ 小 $\to$ 过滤器\emph{激进}（贴近边界才修正，效率高但容错小，模型误差/延迟稍大就可能越界）。这个权衡没有免费午餐，但混合架构提供了一个缓解思路：\emph{让 GCBF+ 这类学习式 CBF 自动找到比手工 CBF 更不保守的安全边界}——这正是 GCBF+ “比手工 CBF 好 $20\%$”的来源。

\begin{pitfall}[概念误区：对二阶/三阶系统用一阶 CBF，过滤器形同虚设]
错误做法：机器人是力控/加速度控制（二阶动力学），却用一阶 CBF 不等式 $\dot h\ge-\alpha h$ 直接约束控制。后果：$\dot h$ 里根本不含控制 $u$，约束对 $u$ 不起作用，过滤器“通过”了所有控制，碰撞照常发生——而你以为有安全保护。根本原因：没检查 CBF 的相对度。正确做法：先确定相对度（$h$ 对 $u$ 求几次导才出现 $u$），相对度 $\ge2$ 时必须用 HOCBF，构造嵌套函数链 $\psi_0,\psi_1,\dots$ 直到显含 $u$。
\end{pitfall}

\begin{pitfall}[编程陷阱：QP infeasible 时退回名义控制，在最危险时刻关闭保护]
错误做法：CBF-QP 求解失败时 \texttt{return u\_nom}。后果：infeasible 通常恰好发生在多个约束冲突的“危险密集逼近”时刻，此时退回名义控制等于在最该保护的瞬间撤掉保护，几乎必撞。根本原因：把 infeasible 当成“无约束”处理，方向完全反了——infeasible 意味着“找不到安全控制”。正确做法：QP infeasible 时降级到保守安全动作（零加速度、最大减速、原地悬停/急停）。更进一步可用软约束 CBF（加松弛变量 $+$ 大惩罚）保证 QP 始终可行。
\end{pitfall}

\begin{pitfall}[思维陷阱：以为安全过滤可以完全替代上层的安全意识]
错误做法：因为有了 CBF 过滤层，就让 MARL 完全不管安全（奖励里去掉所有碰撞惩罚），认为“反正过滤器会兜底”。后果：MARL 学出极度激进的策略（总把控制顶到极限，反正过滤器会拉回来），导致过滤器持续满负荷工作、频繁触及边界，系统在安全边缘反复横跳，效率极低且对模型误差极敏感。根本原因：安全过滤保证“不越界”，但不保证“行为优雅”。正确做法：上层（MARL）仍应保留\emph{温和的}安全意识（适度碰撞惩罚），过滤器只做“最后一道硬保证”，这样它大部分时间透明（零修正）、只在极端情况下介入。
\end{pitfall}

\begin{pitfall}[概念误区：CBF 保证安全就等于保证任务能完成]
错误做法：既然 CBF-QP 保证安全又追踪标称控制，就以为它既安全又一定能完成任务。后果：对称/密集场景会\emph{死锁}（deadlock/livelock）——名义控制要往前（去目标）、CBF 要往后（避障），合力为零，机器人完全安全但僵在原地永远到不了目标。根本原因：CBF 保证\textbf{安全}（前向不变性）但不保证\textbf{活性}（liveness，好事会发生），二者是独立性质。正确做法：把 CBF 当\emph{安全滤波器}而非\emph{完整规划器}；死锁要专门处理——用上层（MARL/规划）提供“谁让谁”的决策打破对称（这正是 \cref{sec:mh-hierarchy} 让 MARL 做高层协调的价值之一），或在 CBF 设计中引入旋转分量/优先级。
\end{pitfall}

\begin{practice}
\begin{enumerate}
  \item \textbf{【实现题】} 基于 \verb@pairwise_cbf_filter@，为 $3$ 个 2D 双积分器实现完整碰撞避免仿真：$3$ 个 agent 初始在等边三角形顶点，名义控制让每个直奔对面顶点（必然交叉于中心）。验证：(a) 任意时刻任意两 agent 距离 $\ge2r$；(b) 画出 $\min_{i\ne j}\|p_i-p_j\|$ 随时间的曲线，确认它始终 $\ge2r$ 且在交汇时刻“贴着”$2r$。再做对比实验：去掉 CBF 过滤，确认 agent 在中心相撞。
  \item \textbf{【调试挑战】} 现象“agent 总是冲到离障碍极近才急刹，偶尔来不及就撞了”，已知代码逻辑结构正确。判断最可能的 bug 在哪（提示：$\alpha$ 的取值），并说明 $\alpha_1,\alpha_2$ 应往哪个方向调、调整背后的物理含义（$\alpha$ 如何控制“多早开始减速”）。
  \item \textbf{【设计扩展 $+$ 跨章综合】} 综合 \cref{sec:mh-hierarchy}、\cref{sec:mh-safety} 与规划部的 MAPF，为“$10$ 台 AGV 在仓库执行终身取送货”设计两级安全方案：(a) 离散级——用 MAPF/优先级规划做无碰撞的全局路径；(b) 连续级——用 CBF 过滤处理“离散计划无法覆盖的连续执行误差和动态障碍”。说明两级各自保证什么、为什么需要两级（提示：离散 MAPF 假设“同步、单位步、点机器人”，真机违反这些假设时谁来兜底），以及 MARL 在这套方案里的角色。
\end{enumerate}
\end{practice}

% ════════════════════════════════════════════════════════════════
\section{残差强化学习：在可信基线上学修正}
\label{sec:mh-residual}

\paragraph{动机：从零学 vs 站在巨人肩上。}
\cref{sec:mh-hierarchy,sec:mh-safety} 讲的“纵向分层”和“安全过滤”都把 RL 和模型放在\emph{不同的层/模块}。这一节讲第三种、也是哲学上最优雅的混合形态——\textbf{残差强化学习}（residual RL）：不分层，而是让传统控制器和 RL 在\emph{同一个控制空间里叠加}，控制器出一个可信的基线 $u_{\text{base}}$，RL 只学一个修正量 $u_{\text{res}}$。设想你要训一个 RL 策略让机械臂精确插孔，或让四足在不平地面稳定行走。纯 RL 从零开始——策略初始随机，机械臂会乱挥、四足会立刻摔倒，要经过几百万步试错才能学会基本能力。这有两个致命问题：(1) \emph{样本效率极低}（大量样本浪费在学习“如何不摔倒”这种早被经典控制解决的基础能力上）；(2) \emph{训练初期极度危险}（随机策略在真机上就是灾难）。但是——\textbf{我们明明已经有一个还不错的控制器了！} 四足有成熟的 MPC/步态控制器（能走，只是地形适应性差一点），机械臂有阻抗控制器（能大致对准，只是精细力控差一点）。这些经典控制器解决了 $90\%$ 的问题，只在最后 $10\%$（未建模的接触、复杂地形、精细力控）上力不从心。\textbf{残差 RL 的核心洞察就是：与其让 RL 从零学那 $100\%$，不如让它站在经典控制器的肩膀上，只学那缺失的 $10\%$。}

\begin{insight}[残差 RL 的工程哲学：不丢弃可信知识，在它之上增量学习]
残差 RL 体现了一种深刻的工程哲学——\textbf{不要丢弃已有的、可信的知识，而要在它之上增量学习}。经典控制器是几十年控制论积累的“先验知识”，它虽不完美但极其可靠。从零训 RL 等于假装这些知识不存在、让神经网络从随机权重重新发现它们——这是巨大的浪费和风险。这和 \cref{sec:mh-decoupling} 解耦原理一脉相承，只是这次的切分维度是“已知 vs 未知”“可靠基础 vs 困难修正”，且切分发生在\emph{同一个控制空间的叠加}上，而非不同的层。
\end{insight}

\paragraph{反面：从零训练相比残差差在哪。}
\begin{itemize}
  \item \textbf{差距一：样本效率差一个数量级。} 残差 RL 的探索围绕一个\emph{已经不错的基线}展开——策略初始残差为零（或接近零），系统行为就是基线控制器的行为（已能完成大部分任务），RL 只需在这个“good enough”起点附近做小幅探索。而从零训练的探索在整个动作空间里盲目游荡，绝大部分撞在“完全不可行”的区域。Johannink 等\cite{johannink2019residual}（ICRA 2019\pz{源稿此处误作 2017；该工作 arXiv 1812.03201 发于 2018-12、正式发表于 ICRA 2019，已据 DOI 10.1109/ICRA.2019.8794127 校正，与下文及 \cref{sec:mh-residual} 统一。}）的实验显示残差 RL 在多个操作任务上样本效率比从零训提升数倍到一个数量级。\pz{“数倍到一个数量级”为源稿转述，具体倍数待核 \cite{johannink2019residual}。}
  \item \textbf{差距二：训练全程的安全性。} 残差 RL 通常对残差做\emph{幅值限制}（$\|u_{\text{res}}\|\le\rho$）——无论 RL 学成什么样，最终控制都不会偏离基线太远；如果基线安全，那么“基线 $+$ 有界残差”也大致安全。这给了残差 RL 一个\emph{安全性下界}：最坏情况不会比基线差太多。而从零训练没有这个保护。
  \item \textbf{差距三：可解释性与可调试性。} 残差 RL 的输出可分解为“基线贡献 $+$ 学习贡献”，你能清楚看到 RL 在每个状态修正了多少、往哪个方向修正，比纯黑盒策略可解释得多。
\end{itemize}

\paragraph{历史：残差学习的思想脉络。}
“残差”这个思想在机器学习里有深厚根基。最著名的是 \textbf{ResNet（残差网络，2015）}\cite{he2016resnet}——它让网络层学习“相对于恒等映射的残差” $F(x)=H(x)-x$ 而非直接学 $H(x)$，因为学“偏离恒等的修正”比学“完整映射”容易得多（当最优映射接近恒等时，残差接近零，极易学）。残差 RL\cite{johannink2019residual}（Johannink 等，ICRA 2019% ；Silver 等同期\cite{silver2005cooperative} —— 误引：silver2005cooperative 为 David Silver 多机寻路论文(AIIDE 2005)，非残差 RL；同期残差工作疑指 Tom Silver《Residual Policy Learning》(2018, arXiv:1812.06298)，待补正确键后恢复
\rebuilt{“Silver 等同期”引用已注释，原键张冠李戴，勿径直重建}）把这个思想搬到控制：让 RL 学“相对于基线控制器的修正”而非完整控制策略，其数学动机和 ResNet 完全一致——当基线已接近最优时，最优残差接近零，RL 极易学到。这条线后来不断深化：\textbf{基线从简单走向复杂}——早期残差 RL 的基线是 P 控制器或手工策略（简单但弱），Residual MPC\cite{residualmpc2025}（2025）把基线升级为 GPU 并行的 MPC（强大且实时），RL 只补 MPC 模型误差导致的那点偏差。主线是：\textbf{基线越强，RL 要学的残差越小、越容易、越安全}。

\begin{insight}[残差 RL 和 ResNet 的共性：最优解接近一个已知的好猜测]
残差 RL 和 ResNet 的深层共性，是都在利用“\textbf{最优解接近一个已知的好猜测}”这个结构。ResNet 假设最优层映射接近恒等，残差 RL 假设最优控制接近基线。在这个假设成立时，“学残差”把一个困难的学习问题（学完整映射/策略）转化成一个容易的学习问题（学一个接近零的小修正）。这也揭示了残差 RL 的适用边界——\textbf{当基线很差（离最优很远）时，残差就不“小”了，残差 RL 的优势会消失}。所以残差 RL 的前提是“你已经有一个还不错的基线”。
\end{insight}

\subsection{理论：残差 RL 的公式化与安全性下界}
\label{sec:mh-residual-theory}

残差 RL 的核心公式简洁到只有一行：
\begin{equation}
  u=u_{\text{base}}(x)+u_{\text{res}}^{\theta}(x),
  \label{eq:mh-residual}
\end{equation}
其中 $u_{\text{base}}(x)$ 是固定的基线控制器（PD、阻抗控制、MPC 等，\textbf{不参与训练}），$u_{\text{res}}^{\theta}(x)$ 是参数为 $\theta$ 的 RL 策略输出的残差。RL 通过标准策略梯度（PPO/SAC 等\cite{levine2018reinforcement}）优化 $\theta$，奖励是任务奖励。虽然公式简单，但有几个关键设计点决定成败：
\begin{description}
  \item[设计点一：残差的初始化与幅值约束。] 训练初期残差应接近零（让系统从“纯基线行为”出发），通常通过\emph{初始化策略网络最后一层权重为接近零}实现；同时对残差做幅值限制 $u_{\text{res}}\in[-\rho,\rho]$（通过 \verb@tanh@ 激活 $\times$ 缩放）。这个上限 $\rho$ 控制“RL 能在多大程度上偏离基线”——上限越小越安全（越贴近基线）但修正能力越弱，上限越大修正能力越强但安全性下界越松。这是残差 RL 最重要的超参数。
  \item[设计点二：残差作用的空间。] 残差可加在不同的控制语义上——加在力矩上（最底层）、加在 MPC 的参考上（中层）、加在目标位置上（高层）。加得越底层修正越直接但越需承受动力学 gap，加得越高层越抽象但修正越间接（回到 \cref{sec:mh-interface} 接口设计的逻辑）。
  \item[设计点三：基线是否随状态切换。] 高级用法里基线本身可以是状态相关的（不同地形用不同步态基线），RL 残差在所有基线之上统一学修正。
\end{description}

\begin{proposition}[残差 RL 的安全性下界]\label{prop:mh-residual-bound}
设基线控制器在所有状态下满足某个安全性质 $S$（如“输出有界、不会发散”），残差幅值限制为 $\|u_{\text{res}}\|\le\rho$。则总控制 $u=u_{\text{base}}+u_{\text{res}}$ 满足“偏离基线不超过 $\rho$”。如果基线的安全裕度大于 $\rho$（即基线“离危险还有 $\rho$ 以上的余量”），则残差 RL 全程不会越过安全边界。
\end{proposition}

\noindent\cref{prop:mh-residual-bound} 是一个\emph{可量化的安全下界}——它把“RL 的不可控性”限制在一个由 $\rho$ 控制的小邻域内，这正是残差 RL 比纯 RL 安全的数学根据。把这个下界和 $\rho$ 的选择联系起来看一个具体直觉：假设基线 MPC 输出的力矩总在执行器极限 $\tau_{\max}=10$ N$\cdot$m 内，且留有 $2$ N$\cdot$m 的裕度（即基线最多用到 $8$ N$\cdot$m）。若把残差限幅设为 $\rho=2$ N$\cdot$m，则总力矩 $\le8+2=10=\tau_{\max}$——残差恰好“吃掉”裕度但不超限，安全下界成立。若贪心地把 $\rho$ 设成 $5$ N$\cdot$m，则总力矩可能到 $8+5=13>\tau_{\max}$，残差能把系统推出执行器极限（饱和、失稳）——安全下界被打破。这就是为什么 $\rho$ 必须\textbf{按基线的实际裕度来定}，而非越大越好：$\rho$ 是“借给 RL 的自由度额度”，借出去的不能超过基线攒下的安全余量。

\begin{insight}[保守的基线和强力的残差相互成全]
\cref{prop:mh-residual-bound} 解释了一个反直觉的现象——\textbf{基线越保守（裕度越大），能安全分配给残差的额度 $\rho$ 反而越大，RL 的修正空间越大}。保守的基线和强力的残差，在安全下界这个约束下是相互成全的：你不必在“保守基线”和“强力学习”之间二选一，恰恰是保守的基线攒下的余量，才让你能放心地给 RL 更大的探索自由度。
\end{insight}

\subsection{理论-工程桥接：在 PD/MPC 基线上叠加残差策略}
\label{sec:mh-residual-code}

把残差 RL 的核心结构落地。两个关键：(i) 残差策略\emph{最后一层}权重初始化为接近零（\verb@uniform(-1e-3,1e-3)@）、偏置为零——这样训练第一步残差 $\approx0$、系统从纯基线行为出发；(ii) 用 \verb@tanh@ $\times$ \verb@res_max@ 把残差锁在 $[-\rho,\rho]$ 内（\cref{prop:mh-residual-bound} 的安全下界）。基线 \verb@baseline_fn@ 可以是 PD、阻抗控制或 \verb@MPC.solve@，它\emph{不参与训练}（不放进优化器）。

\begin{codebox}[Python]{残差控制器：$u=u_{\text{base}}+u_{\text{res}}^{\theta}$（末层零初始化 $+$ tanh 限幅）}
import torch, torch.nn as nn, numpy as np

class ResidualPolicy(nn.Module):
    def __init__(self, obs_dim, act_dim, res_max, hidden=256):
        super().__init__()
        self.res_max = res_max                  # rho: amplitude bound (safety LB)
        self.net = nn.Sequential(
            nn.Linear(obs_dim, hidden), nn.ELU(),
            nn.Linear(hidden, hidden), nn.ELU(),
            nn.Linear(hidden, act_dim))
        nn.init.uniform_(self.net[-1].weight, -1e-3, 1e-3)   # init residual ~ 0
        nn.init.zeros_(self.net[-1].bias)

    def forward(self, obs):
        return self.res_max * torch.tanh(self.net(obs))      # bounded in [-rho, rho]

class ResidualController:                       # u = u_base(x) + u_res^theta(x)
    def __init__(self, baseline_fn, residual_policy):
        self.baseline = baseline_fn             # FROZEN (not in optimizer)
        self.policy = residual_policy

    def act(self, obs, state):
        u_base = self.baseline(state)           # reliable 90% (NOT trained)
        with torch.no_grad():
            u_res = self.policy(torch.as_tensor(obs, dtype=torch.float32)).numpy()
        return u_base + u_res, u_base, u_res    # return decomposition for debugging
\end{codebox}

\noindent\textbf{四处易错}：(i) 残差策略末层\emph{随机}初始化 $\to$ 第一步残差很大、把系统踹到混乱状态，退化成“带偏置的从零训练”；必须末层零初始化。(ii) 残差\emph{不限幅} $\to$ 偏离基线任意远、突破基线裕度，退化成不安全的纯 RL；必须 \verb@tanh@$\times\rho$。(iii) 把基线参数也放进优化器 $\to$ 基线被 RL 梯度往“高奖励但可能不安全”的方向带偏，锚点漂移、安全论证失效；基线必须\emph{冻结}。(iv) 基线太弱（如恒为零）$\to$ $u=0+u_{\text{res}}$ 完全退化成从零训练，残差不“小”、优势荡然无存。

\subsection{残差 RL 在多机器人场景的特殊价值}
\label{sec:mh-residual-multibot}

前面以单机为主，但残差 RL 在\textbf{多机器人}场景有额外的独特价值——这正是本章关心的。
\textbf{价值一：多体交互的“残差”恰是难建模的部分。} 单机控制器（MPC）对单个机器人的动力学建模很准，但它\emph{看不见、也难以建模其他机器人对自己的真实影响}（气流耦合、协同搬运时的力交互、密集编队的尾流）。这些“多体交互残差”恰是单机 MPC 的盲区，而 MARL 残差能从多机交互数据中学到它。于是结构变成：每台机器人用单机 MPC 作基线（搞定自身动力学），用一个 MARL 残差策略补“多体交互”那部分。
\textbf{价值二：残差天然继承单机控制器的成熟度。} 多机器人系统的单机控制（足式 MPC、无人机几何控制）往往已非常成熟，从零训一个多机策略等于把这些成熟成果全扔掉重学。残差 RL 让你复用每台机器人成熟的单机控制器，只在多机协同层面增量学习。
\textbf{价值三：残差限幅提供多机安全的一层保障。} 多机场景一个 agent 的异常会传染，残差限幅保证每台机器人都不会偏离其单机基线太远，从而限制了“单点异常”的幅度，配合 \cref{sec:mh-safety} 的 CBF 过滤形成双重保障。

\begin{insight}[三种混合形态可沿不同维度同时切分、叠加]
在多机器人场景，残差 RL 和分层混合其实可以“合体”——把分层混合底层的单机 MPC 当作残差 RL 的基线，让 MARL 同时承担两个角色：在高层出协调命令（分层混合的 RL 角色），在底层出多体交互残差（残差 RL 的 RL 角色）。这揭示一个重要事实：\textbf{三种混合形态不是非此即彼的菜单，而是可以沿不同维度同时切分、叠加使用的正交技巧}。一个成熟的多机器人混合系统，往往同时是“分层的 $+$ 残差的 $+$ 带安全过滤的”——这正是 \cref{sec:mh-safemarl} 选型框架要帮你组织的。
\end{insight}

\begin{pitfall}[编程陷阱：残差策略最后一层随机初始化，第一步就破坏基线]
错误做法：残差策略网络用默认随机初始化，训练初期残差幅值很大。后果：训练第一步系统行为就完全偏离基线（被随机大残差推到混乱状态），残差 RL 退化成“带一个固定偏置的从零训练”，丢掉基线的全部样本效率和安全优势。根本原因：残差 RL 的核心前提是“从基线行为出发、小幅探索”。正确做法：把残差策略网络\textbf{最后一层}权重初始化为接近零（如 \texttt{uniform(-1e-3,1e-3)}）、偏置为零。
\end{pitfall}

\begin{pitfall}[思维陷阱：残差不限幅，退化成不安全的纯 RL]
错误做法：残差策略直接输出无界的修正量，不做幅值限制。后果：残差可无限增长，总控制偏离基线任意远，基线的安全裕度被突破，残差 RL 失去“安全性下界”，和纯 RL 一样不安全。根本原因：残差 RL 的安全性来自“总控制不会偏离可信基线太远”（\cref{prop:mh-residual-bound}），这个“太远”必须用残差幅值上限 $\rho$ 显式约束。正确做法：用 \texttt{tanh}$\times\rho$ 把残差锁在 $[-\rho,\rho]$ 内，$\rho$ 按基线安全裕度来定。
\end{pitfall}

\begin{pitfall}[概念误区：把基线控制器也加进训练，污染其可信性]
错误做法：为“让基线也变得更好”，把基线控制器的参数也放进优化器一起训。后果：基线被 RL 梯度往“高奖励”方向带偏，可能牺牲它原本的安全/稳定保证，残差 RL 退化成纯 RL。根本原因：基线的全部价值在于它“可信、固定、有理论保证”，是残差 RL 的“安全锚点”，一旦让它可训练，锚点就漂移了。正确做法：基线控制器必须\textbf{冻结}（参数不放进优化器，或对其输出 \texttt{detach}）；若确实想改进基线，应在 RL 之外用控制论方法单独改进，改完再冻结。
\end{pitfall}

\begin{practice}
\begin{enumerate}
  \item \textbf{【实现题】} 实现一个完整的残差 RL 训练循环（PPO 或 SAC），任务是“2D 点质量精确到达目标”，基线是一个\emph{故意调得偏弱}的 PD 控制器（增益偏低，会有稳态误差和缓慢响应）。验证：(a) 训练初期（残差 $\approx0$）系统行为就是 PD 基线行为；(b) 训练后残差学会补偿 PD 的稳态误差和响应迟缓；(c) 画出“基线贡献 $u_{\text{base}}$”和“残差贡献 $u_{\text{res}}$”随训练进程的幅值变化。
  \item \textbf{【对比实验】} 用同一任务和同一 RL 算法，对比“从零训练”（基线恒为零）vs “残差 RL”（PD 基线）的样本效率（画“奖励 vs 训练步数”曲线）。然后做关键消融——把残差 RL 的基线从“还不错的 PD”换成“很差的 PD”（增益极低），观察残差 RL 的优势如何随基线质量下降而消失（亲手验证 \cref{prop:mh-residual-bound} 背后的洞察）。
  \item \textbf{【设计扩展 $+$ 跨章综合】} 综合 \cref{sec:mh-interface,sec:mh-residual} 与控制部 MPC，为“$4$ 架无人机密集编队穿越障碍”设计一个\textbf{分层 $+$ 残差}的混合控制器：(a) 每架无人机用单机几何控制/MPC 作基线；(b) 一个 MARL 策略同时输出“高层编队命令”（\cref{sec:mh-interface}）和“底层多体交互残差”（\cref{sec:mh-residual}，补气流耦合）。画出数据流图，明确 MARL 的两个输出分别加在哪里、为什么“多体交互”适合用残差而“编队决策”适合用分层命令。再说明你会如何用 \cref{sec:mh-safety} 的 CBF 过滤给这个系统兜底硬安全。
\end{enumerate}
\end{practice}

% ════════════════════════════════════════════════════════════════
\section{Safe MARL 与三范式对比收口}
\label{sec:mh-safemarl}

\paragraph{动机：安全能不能“训”进策略里？}
\cref{sec:mh-safety} 的安全过滤是一种“\emph{外挂式}”安全——策略该怎么训怎么训，部署时在外面套一个 CBF 过滤器兜底。这很有效，但有一个微妙的不足：\textbf{策略本身从未“意识到”安全约束的存在}，它学到的是“在没有约束的世界里最优”的策略，过滤器在外面“硬掰”它的输出。这会导致策略和过滤器持续“对抗”（策略想越界、过滤器拉回来），既损失性能又把系统推到安全边缘。一个互补的思路是：\textbf{能不能在训练时就把安全约束告诉策略，让它学会“在约束下最优”？} 这样策略本身就“知道”该避让，过滤器（如果还用的话）大部分时间无需介入。这就是 Safe MARL 的核心——把安全写进\emph{优化目标的约束}，而不是事后过滤。注意这和 \cref{sec:mh-why} “安全必须是约束而非目标项”的洞察一脉相承——只不过这次的“约束”作用在\textbf{训练优化}上（让策略学会满足期望意义的约束），而 \cref{sec:mh-safety} 的“约束”作用在\textbf{逐时刻执行}上（硬保证）。两者是安全的两个层次，互补而非替代。

\paragraph{反面：如果安全只靠“奖励惩罚”会怎样。}
\cref{sec:mh-why} 已论证纯奖励惩罚的根本缺陷。这里从“约束优化”视角再深化一层。把安全写成奖励里的惩罚项 $r=r_{\text{task}}-\lambda\cdot c_{\text{safety}}$，$\lambda$ 是手工权重。问题在于 $\lambda$ 是\emph{固定的、手调的}：$\lambda$ 太小 $\to$ 策略为任务奖励牺牲安全（碰撞频繁）；$\lambda$ 太大 $\to$ 策略过度保守且训练不稳定；\textbf{更糟的是没有“正确的 $\lambda$”}——不同任务阶段、不同状态需要的安全-性能权衡不同，一个固定 $\lambda$ 无法适应。约束公式化的精妙在于：\textbf{它让 $\lambda$（现在叫 Lagrangian 乘子 $\mu$）自动、自适应地调整}——当策略违反安全约束时 $\mu$ 自动增大（更重视安全），当策略安全裕度充足时 $\mu$ 自动减小（更重视性能）。这把“手调一个固定权重”变成“自动求解一个约束优化”。

\paragraph{历史：从 CMDP 到约束马尔可夫博弈。}
约束强化学习的理论根基是\textbf{约束马尔可夫决策过程}（Constrained MDP, CMDP）\cite{altman1999cmdp}（Altman, 1999）——在标准 MDP 上加一个“累积代价的期望 $\le$ 阈值”的约束。求解 CMDP 的主流方法是 \textbf{Lagrangian 松弛}（把约束用乘子吸进目标）和 \textbf{CPO}（Constrained Policy Optimization, 2017\cite{achiam2017constrained}，在信赖域内保证每步更新都满足约束）。把 CMDP 推广到多智能体，就得到\textbf{约束马尔可夫博弈}（Constrained Markov Game, CMG）——每个 agent 既要最大化自己的累积奖励，又要满足自己的累积安全代价约束。把任务公式化为 CMG、用 \textbf{HAPPO $+$ Lagrangian} 求解、并对安全代价做\emph{分解}（碰撞 $+$ 翻转 $+$ 关节限位违反分别约束）的双四足协同搬运工作是这条线的代表\rebuilt{（源稿归为 Jose \& Zhang 2025，归属/年份待联网核对）}。这条线和 GCBF+（安全过滤一侧）、residual RL（残差一侧）在 2024--2025 年汇成了“Safe MARL”这个活跃方向（HAPPO 见 \cref{ch:plan-safemarl}\cite{kuba2022happo}）。

\subsection{理论：约束马尔可夫博弈与 Lagrangian 松弛}
\label{sec:mh-cmg}

正式写出 CMG 并推导 Lagrangian 求解。回顾标准 MARL 的目标——每个 agent $i$ 最大化期望累积折扣奖励 $J_i^{\text{reward}}(\theta)=\mathbb{E}[\sum_t\gamma^tr_{i,t}]$。CMG 在此基础上加一个\textbf{安全约束}——每个 agent 的期望累积安全代价不超过阈值 $\delta$：
\begin{equation}
  \max_{\theta_i}\ J_i^{\text{reward}}(\theta)=\mathbb{E}\Big[\textstyle\sum_t\gamma^tr_{i,t}\Big]\quad\text{s.t.}\quad J_i^{\text{cost}}(\theta)=\mathbb{E}\Big[\textstyle\sum_t\gamma^tc_{i,t}\Big]\le\delta,
  \label{eq:mh-cmg}
\end{equation}
其中安全代价 $c_{i,t}$ 是可分解的（工程上的关键设计）：$c_{i,t}=c_{i,t}^{\text{collision}}+c_{i,t}^{\text{flip}}+c_{i,t}^{\text{joint-limit}}+\dots$。约束 $J^{\text{cost}}\le\delta$ 直接表达了“我能容忍多少不安全”（$\delta$ 有清晰物理含义，如“期望碰撞代价不超过 $0.01$”），而惩罚权重 $\lambda$ 没有这种直接含义——$\delta$ 是可解释、可设定的，$\lambda$ 是需要试凑的。

\begin{theorem}[CMG 的 Lagrangian 原始-对偶更新]\label{thm:mh-lagrangian}
把约束用乘子 $\mu\ge0$ 吸进目标，构造 Lagrangian $\mathcal{L}(\theta,\mu)=J^{\text{reward}}(\theta)-\mu\big(J^{\text{cost}}(\theta)-\delta\big)$，得到一个 min-max（鞍点）问题 $\max_\theta\min_{\mu\ge0}\mathcal{L}$。用原始-对偶交替优化求解：
\begin{align}
  \theta &\leftarrow \theta+\eta_\theta\,\nabla_\theta\big[A^{\text{reward}}-\mu\,A^{\text{cost}}\big] &&\text{(策略更新：HAPPO step)}\label{eq:mh-pd-theta}\\
  \mu &\leftarrow \max\!\big(0,\ \mu+\eta_\mu\,(J^{\text{cost}}-\delta)\big) &&\text{(乘子更新：约束违反则增大)}\label{eq:mh-pd-mu}
\end{align}
其中 $A^{\text{reward}},A^{\text{cost}}$ 分别是奖励/代价优势函数。收敛要求 $\eta_\mu\ll\eta_\theta$（对偶变量更新比原始变量慢）。
\end{theorem}
\begin{proof}[逐行理解]\rebuilt{（标准 CMDP 原始-对偶结果\cite{altman1999cmdp}，多智能体推广用 HAPPO\cite{kuba2022happo}）}
\textbf{策略更新 \cref{eq:mh-pd-theta}}：用 HAPPO（异构 agent 序贯更新，\cref{ch:plan-safemarl}）优化一个\emph{组合优势} $A^{\text{reward}}-\mu A^{\text{cost}}$——既追求奖励优势，又（被 $\mu$ 加权地）压低代价优势；$\mu$ 越大，策略越重视降低安全代价。\textbf{乘子更新 \cref{eq:mh-pd-mu}}：$\mu$ 朝“约束违反量” $J^{\text{cost}}-\delta$ 的方向走——若当前策略\emph{违反}约束（$J^{\text{cost}}>\delta$），$\mu$ 增大 $\to$ 下一步策略更重视安全；若\emph{满足}约束（$J^{\text{cost}}<\delta$），$\mu$ 减小 $\to$ 下一步可更重视性能；$\max(0,\cdot)$ 保证 $\mu\ge0$（不等式约束的乘子非负）。这正是投影次梯度上升在对偶变量上的应用。
\end{proof}

\begin{insight}[Lagrangian 乘子 $\mu$ 是一个自动的安全-性能权衡控制器]
$\mu$ 的自适应更新本质上是一个\textbf{自动的安全-性能权衡控制器}——它实时测量“当前策略离安全约束有多远”，据此动态调整对安全的重视程度。这正是它优于固定惩罚权重 $\lambda$ 的根本原因：固定 $\lambda$ 是“开环”地设定权衡，$\mu$ 是“闭环”地调节权衡（用约束违反量做反馈）。把这个洞察和控制论联系起来——$\mu$ 的更新 \cref{eq:mh-pd-mu} 就是一个\textbf{积分控制器}：它积累“约束违反量”并据此调节，直到约束恰好满足（违反量为零、$\mu$ 稳定）。安全约束优化和反馈控制在这里殊途同归。这也解释了 \cref{thm:mh-lagrangian} 为何要求 $\eta_\mu\ll\eta_\theta$——又一次时间尺度分离：慢变量（$\mu$）调节快变量（$\theta$）的边界条件。
\end{insight}

\subsection{两种安全语义的根本差异}
\label{sec:mh-two-semantics}

现在我们能清晰对比本章出现的两种安全机制了，它们有\textbf{根本不同的安全语义}（\cref{tab:mh-two-safety}）。

\begin{table}[H]\centering\small
\caption{两种安全语义：Safe MARL（CMG）vs 安全过滤（CBF）}
\label{tab:mh-two-safety}
\begin{tabular}{p{3.0cm}p{4.6cm}p{4.4cm}}
\toprule
 & Safe MARL（CMG $+$ Lagrangian） & 安全过滤（CBF-QP / GCBF+） \\
\midrule
安全语义 & 期望意义（$\mathbb{E}[\text{cost}]\le\delta$） & 逐状态硬保证（每帧 $h\ge0$） \\
作用时机 & 训练时（策略学会约束） & 部署时（执行前过滤） \\
保证强度 & 软（期望满足，单次仍可能违反） & 硬（前向不变性，永不违反） \\
是否需运行时模块 & 否（约束已内化进策略） & 是（部署时 QP 必须在环） \\
对模型的依赖 & 不需要动力学模型 & 需要 $f,g$（控制仿射模型） \\
典型适用 & 安全偏好（少碰撞但非零容忍） & 安全关键（绝对不可碰撞） \\
\bottomrule
\end{tabular}
\end{table}

\begin{insight}[Safe MARL 和安全过滤是安全的两道防线]
Safe MARL 和安全过滤不是竞争关系，而是\textbf{安全的两道防线}。Safe MARL 让策略“通常很安全”（期望意义，训练内化），把安全代价压到很低——这减轻了过滤器的负担；安全过滤做“绝对兜底”（逐状态硬保证，部署强制）——它接住 Safe MARL 漏掉的极端情况。最强的安全架构是\textbf{两者叠加}：用 Safe MARL 训出一个安全裕度很大的策略（让过滤器大部分时间透明），再用 CBF 过滤兜底（保证万无一失）。这呼应了 \cref{sec:mh-residual} 末尾的洞察——成熟系统总是多种安全机制叠加，因为没有单一机制能同时做到“高性能 $+$ 期望安全 $+$ 硬保证”。
\end{insight}

\subsection{三大范式对比与选型框架（多机器人方向收口）}
\label{sec:mh-selection}

现在站到最高处，把整个多机器人方向学过的东西收口。面对一个真实的多机器人任务，你有三大范式可选——纯 MARL、纯传统规控、混合。\cref{tab:mh-paradigm} 是它们的“全身像”。

\begin{table}[H]\centering\small
\caption{三大范式对比（多机器人方向收口）。$\checkmark$ 强 / $\times$ 弱 / $\circ$ 中}
\label{tab:mh-paradigm}
\begin{tabular}{p{2.6cm}p{3.0cm}p{3.2cm}p{4.0cm}}
\toprule
维度 & 纯 MARL & 纯传统规控 & 混合（分层/过滤/残差） \\
\midrule
安全保证 & $\times$ 软（概率非零） & $\checkmark$ 硬（CBF/MPC 约束） & $\checkmark$ 硬（CBF 过滤兜底） \\
难建模决策 & $\checkmark$ 强（自动学习） & $\times$ 弱（手工代价/规则） & $\checkmark$ 强（RL 出高层） \\
样本效率 & $\times$ 低（从零试错） & --- 无需训练 & $\circ$ 中-高（残差/分层缩小学习面） \\
可解释/可调试 & $\times$ 黑盒 & $\checkmark$ 每步可查 & $\circ$ 中（底层可查，高层黑盒） \\
sim-to-real & $\times$ 难（gap 弥漫全链） & $\checkmark$ 易（用真实模型） & $\checkmark$ 较易（gap 局部化到高层接口） \\
可扩展性($N$) & $\circ$ 中 & $\times$ 差（$N^2$ 约束/联合爆炸） & $\checkmark$ 好（GCBF+/分布式 MPC） \\
工程复杂度 & $\circ$ 中（训练管线） & $\circ$ 中（建模 $+$ 调参） & $\times$ 高（多模块 $+$ 接口） \\
典型场景 & 纯仿真、研究、难建模博弈 & 模型准确、单一任务、安全关键 & 真实多机器人系统（绝大多数） \\
\bottomrule
\end{tabular}
\end{table}

\noindent\textbf{读这张表的关键结论}：对绝大多数真实多机器人任务，混合是默认答案——因为真实任务几乎总是同时需要“硬安全保证 $+$ 难建模决策能力 $+$ sim-to-real 可迁移性 $+$ 可扩展性”，而只有混合能同时给出这四样。纯 MARL 和纯规控各自缺了其中一两样。混合的唯一代价是工程复杂度高（多模块、多接口）——这是为“兼得各家之长”付的工程税。

\begin{insight}[“混合是默认”是被两种范式的数学边界逼出来的]
“混合是默认”不是一句口号，而是由两种范式各自的数学边界\emph{逼}出来的必然。用一句反事实锚定它：\textbf{如果某团队声称用“纯端到端 MARL”解决了一个真实部署的安全关键多机器人任务、且没有任何模型化的安全层，那么大概率要么任务其实不安全关键（碰一下没关系），要么他们在某处偷偷藏了一个传统规控的兜底而没明说}——因为纯 MARL 在数学上无法提供硬安全保证，这不是工程努力能弥补的。反过来，如果一个团队用“纯传统规控”做出了灵活的多机协调，那大概率是因为他们的“协调”其实很规整（能写成清晰的优化/规则），一旦协调决策真正进入“难以解析建模”的开放领域，纯规控就会撞墙。这也是本章把 \cref{sec:mh-why} 的“两难”放在最前面讲的原因：理解了边界，选型就成了顺理成章的事。
\end{insight}

\paragraph{选型决策框架。}
沿四个问题逐层判断，给你一个可操作的决策树：
\begin{enumerate}
  \item \textbf{任务是否要求“硬安全保证”}（一次违反就是灾难，如碰撞 $=$ 硬件损毁）？\emph{是} $\to$ 必须有安全层，排除“纯 MARL”，继续问 2；\emph{否}（软安全偏好即可，如游戏/仿真）$\to$ 纯 MARL 可考虑，继续问 2。
  \item \textbf{是否存在“难以解析建模的高层协调决策”}（让行/编队重组/任务再分配）？\emph{是} $\to$ 需要学习层，排除“纯传统规控”，倾向\textbf{分层混合}（RL 出高层）；\emph{否}（决策能写成清晰优化/规则）$\to$ 纯传统规控可能够用，继续问 3。
  \item \textbf{是否已有一个“还不错的单机/基线控制器”、只想补它的不足}？\emph{是} $\to$ 倾向\textbf{残差 RL}（在基线上学修正）；\emph{否} $\to$ 看问 2 结论。
  \item \textbf{$N$ 是否很大}（几十到上千）\textbf{且需要分布式安全}？\emph{是} $\to$ 安全层用\textbf{GCBF+}（GNN 参数化，可扩展），高层用分布式 MARL；\emph{否}（$N$ 小）$\to$ 安全层用手工 CBF-QP 即可，高层 MAPPO/HAPPO。
\end{enumerate}
\textbf{综合}：真实系统的答案常是“以上皆是”的组合——分层混合（RL 高层 $+$ MPC 底层）$+$ 残差 RL（MPC 上叠多体交互残差）$+$ 安全过滤（GCBF+ 兜底硬安全）$+$ Safe MARL（训练时约束让策略学会安全），各机制沿不同维度切分、叠加，各管一段。

\begin{insight}[全章收口：没有银弹，只有组合]
学完整个多机器人方向，最重要的一个认识是——\textbf{没有银弹，只有组合}。通信图、共识、任务分配、MAPF、编队、分布式 MPC、MARL、CBF……这些不是相互竞争的“流派”，而是一个工具箱里互补的工具。一个真实的多机器人系统，是把这些工具沿“离散/连续”“学习/模型”“高层/底层”“决策/安全”“基础/修正”等多个正交维度切分、组装出来的。本章的混合范式，正是这个“组装”的方法论——它教你的不是某个具体算法，而是\textbf{如何判断每个工具的比较优势、如何沿正确的维度切分问题、如何设计模块间的接口}。这才是从“会用算法”到“会设计系统”的关键一跃。
\end{insight}

\begin{pitfall}[概念误区：把 Safe MARL 的“期望安全”误当成“硬安全保证”]
错误做法：用 CMG $+$ Lagrangian 训出一个满足 $\mathbb{E}[\text{cost}]\le\delta$ 的策略，就认为它“安全了”，直接上安全关键场景、不加任何运行时过滤。后果：期望意义的约束只保证“平均而言安全代价低”，单条轨迹、单个时刻仍可能违反——安全关键场景下偶发碰撞、硬件损毁。根本原因：混淆了期望约束（$\mathbb{E}[\cdot]\le\delta$，软）和逐状态硬约束（每帧 $h\ge0$，硬），Lagrangian 求解的是前者。正确做法：安全关键场景下 Safe MARL 只是“第一道防线”，必须再叠加 \cref{sec:mh-safety} 的 CBF 过滤做“硬兜底”。
\end{pitfall}

\begin{pitfall}[思维陷阱：迷信“混合一定比纯方案好”，过度工程化]
错误做法：不管任务多简单，一律上“分层 $+$ 残差 $+$ CBF $+$ Safe MARL”的全套混合，认为越复杂越先进。后果：对一个本可用纯 MPC 干净解决的简单任务（如单机轨迹跟踪），堆了一大套混合架构，引入大量接口 bug、调试噩梦、维护成本，性能反而不如简单方案。根本原因：没理解混合的代价是“工程复杂度”。正确做法：用选型决策框架诚实评估任务真正需要什么——如果任务模型准确、无难建模决策、$N$ 小，纯传统规控就够了。\textbf{复杂度要花在刀刃上}：只为任务真正需要的能力引入对应的混合机制。
\end{pitfall}

\begin{pitfall}[编程陷阱：Lagrangian 乘子 $\mu$ 学习率与策略学习率不匹配导致振荡]
错误做法：$\mu$ 的更新步长 $\eta_\mu$ 设得和策略学习率同量级或更大。后果：$\mu$ 剧烈震荡（约束一违反就猛增、一满足就猛降），导致策略在“过度保守”和“过度激进”之间反复横跳、训练不收敛。根本原因：原始-对偶优化要求对偶变量（$\mu$）更新比原始变量（$\theta$）\emph{慢}（时间尺度分离，\cref{thm:mh-lagrangian}），让策略有时间在固定 $\mu$ 下收敛、$\mu$ 再缓慢调整。正确做法：$\eta_\mu$ 通常比策略学习率小 $1$--$2$ 个数量级，让 $\mu$ 缓慢平滑地调整（这又是一次时间尺度分离的应用，和 $\mu$ 作为积分控制器的洞察一脉相承）。
\end{pitfall}

\begin{practice}
\begin{enumerate}
  \item \textbf{【推导题】} 从 CMG 的约束优化问题 \cref{eq:mh-cmg} 出发，自己完整推导一遍 Lagrangian 原始-对偶更新。回答：(a) 为什么 $\mu$ 的更新方向是 $J^{\text{cost}}-\delta$（而非 $\delta-J^{\text{cost}}$）？(b) 为什么要 $\max(0,\cdot)$？(c) 把 $\mu$ 的更新和“积分控制器”类比——它积分的是什么量、调节的是什么、稳态（收敛）时 $\mu$ 和约束的关系是什么？
  \item \textbf{【选型实战】} 给定三个真实任务，分别用“选型决策框架”判断该用哪种（组合）范式，并完整说明沿四个问题的判断过程：(a) $1000$ 架无人机的灯光秀编队（密集、需绝对不撞、队形复杂）；(b) 单台机械臂在已知工件上的精密装配（模型准确、有成熟阻抗控制器、只差精细力控）；(c) $5$ 台异构机器人在未知灾后环境的协同搜救（难建模决策多、动态障碍、需安全）。
  \item \textbf{【综合设计 $+$ 全章串联】} 这是本章压轴综合题。为“$8$ 台异构机器人（$4$ 无人机 $+$ $4$ 地面车）在动态仓库协同搬运”（即 \cref{sec:mh-why} 开头的任务）设计一套\textbf{完整的混合系统}，串联本章所有小节：(a) 画出完整的分层架构（\cref{sec:mh-hierarchy}）；(b) 设计 RL$\to$MPC 命令接口（\cref{sec:mh-interface}）；(c) 设计 CBF/GCBF+ 安全过滤层（\cref{sec:mh-safety}）；(d) 指出哪里可以用残差 RL（\cref{sec:mh-residual}）；(e) 用 CMG 写出训练时的安全约束（\cref{sec:mh-safemarl}）；(f) 说明每一类 sim-to-real gap（\cref{tab:mh-gap}）被哪个模块吸收。
\end{enumerate}
\end{practice}

% ════════════════════════════════════════════════════════════════
\section{前沿延伸：LLM 作为高层决策层（RoCo 范式）}
\label{sec:mh-llm}

\paragraph{动机：高层决策器还能更“聪明”吗。}
\cref{sec:mh-hierarchy} 的四层架构里，高层决策器（Level 3）一直是一个 MARL 策略网络。但 2024 年以来一个引人注目的方向是——\textbf{把高层决策器从“策略网络”换成“大语言模型”（Large Language Model, LLM）}。这看起来是个大跳跃，但从本章的解耦原理看，它其实是同一框架的自然延伸：LLM 只是另一种“高层决策器”，底层仍是 MPC/CBF。MARL 高层很强，但有几个固有局限：\emph{需要大量训练}（每个新任务都要重新训练或微调）；\emph{难以利用常识}（它不“懂”任务的语义，不知道“易碎品要轻拿轻放”这类人类常识，除非显式编码进奖励）；\emph{难以接受自然语言指令}；\emph{难以解释决策理由}。而 LLM 恰好在这几点上有惊人的互补优势——内置海量人类常识、能理解自然语言指令、能“说出”决策理由（思维链）。

\begin{insight}[从“MARL 高层”到“LLM 高层”没有改变解耦原理]
这个跳跃本质上没有改变本章的解耦原理——它改变的只是“高层决策器用什么实现”。MARL 高层是“用强化学习从奖励中学出协调策略”，LLM 高层是“用预训练语言模型从常识和指令中推理出协调策略”。两者都是 \cref{fig:mh-stack} 四层架构里的 Level 3，都通过同一个命令接口（\cref{sec:mh-interface}）连接到底层 MPC。这说明本章的架构骨架有强大的\textbf{通用性}——高层决策器是可替换的插件，无论它是策略网络还是 LLM，底层的执行与安全机制（MPC/CBF）保持不变。理解了这一点，你就不会被“LLM $+$ 机器人”的新潮表象迷惑，而能看清它在混合范式里的准确位置。
\end{insight}

\paragraph{反面：如果让 LLM 直接控制机器人底层会怎样。}
有人可能想：LLM 这么强，能不能让它直接输出每个关节的力矩或每一步的速度指令？这是一个灾难性的想法，原因和 \cref{sec:mh-interface} “RL 不应直接出底层动作”如出一辙，但更严重：\emph{频率完全不匹配}（LLM 推理一次要几百毫秒到几秒，根本无法满足底层 $50$--$1000$ Hz）；\emph{没有任何动力学概念}（它生成的“力矩”只是符合语言模式的数字）；\emph{没有安全保证}（LLM 会“幻觉”，可能自信地输出完全错误的指令）；\emph{数值精度差}。所以正确的做法必然是——\textbf{LLM 只做最高层的语义决策（慢、抽象、符号化），把所有连续控制留给底层 MPC}。

\paragraph{历史：从符号规划到 LLM 协商。}
让“高层符号决策 $+$ 底层连续控制”分层，这个思想在机器人学里很古老——经典的 \textbf{TAMP}（Task and Motion Planning）就是“高层用符号规划器（如 PDDL）决定动作序列，底层用运动规划器执行”。但经典 TAMP 的高层符号规划器是\emph{人工设计规则 $+$ 搜索}，难以处理开放、模糊、需要常识的任务。LLM 的出现给这个古老分层注入了新活力——\textbf{用 LLM 替换符号规划器}。SayCan\cite{ahn2022saycan}（2022）是早期代表（LLM 提议动作、用价值函数评估可行性）；\textbf{RoCo}\cite{mandi2024roco}（Mandi, Jain, Song, ICRA 2024）则把它推向\emph{多机器人}：多个 LLM agent 通过“辩证对话”（dialectic）协商任务分配和动作，再交给底层 MPC 执行。

\subsection{理论：RoCo 的三层架构}
\label{sec:mh-roco}

RoCo 的架构本质上是本章四层架构的一个变体——把 Level 3（MARL）换成了“LLM 决策层 $+$ LLM 协商层”两个子层，底层仍是 MPC（\cref{fig:mh-roco}）。

\begin{figure}[H]
\centering
\begin{tikzpicture}[
  >=Stealth, font=\small,
  bx/.style={draw, rounded corners, align=left, inner sep=5pt, text width=9.6cm},
  d/.style={bx, fill=second!10, draw=second!60!black},
  n/.style={bx, fill=third!12, draw=third!60!black},
  c/.style={bx, fill=main!8, draw=main!60!black},
  ln/.style={->, thick, gray!55!black}]
\node[d] (dec) {\textbf{LLM 决策层}（每个机器人一个 LLM agent）\\[2pt]
  \footnotesize 输入：自然语言任务 $+$ 文本化场景；输出：每机的子目标序列（符号化，如“抓取红块”）。利用 LLM 的常识与语言理解。};
\node[n, below=6mm of dec] (neg) {\textbf{LLM 协商层（Dialectic，多 agent 对话）}\\[2pt]
  \footnotesize 多个 LLM agent 通过多轮自然语言对话发现并解决冲突（“我要先过通道”——“不行，我的货更急，你等等”）。输出：协商一致、无冲突的联合子目标。\emph{用“语言对话”实现难建模的高层协调决策}。};
\node[c, below=6mm of neg] (ctl) {\textbf{控制层（MPC，每个机器人独立）}\\[2pt]
  \footnotesize 输入：协商后的几何子目标；输出：连续轨迹/控制，到达子目标 $+$ 局部避障。方法：标准 MPC（\cref{fig:mh-stack} 的 Level 2）。};
\draw[ln] (dec)--(neg);
\draw[ln] (neg)-- node[right, font=\scriptsize]{子目标（符号 $\to$ 几何，需 grounding）} (ctl);
\end{tikzpicture}
\caption{RoCo 的三层架构（本章四层架构的变体）：Level 3 被拆成“LLM 决策层 $+$ LLM 协商层”，底层仍是 MPC。注意它缺少高频 WBC 层——部署到真实力控机器人需补上。}
\label{fig:mh-roco}
\end{figure}

\noindent\textbf{逐层解读}：\emph{LLM 决策层}——每个机器人配一个 LLM agent，读取文本化的任务和场景、输出自己的子目标计划，利用 LLM 的常识“知道”任务该怎么分解；\emph{LLM 协商层（Dialectic）}——这是 RoCo 的核心创新，多个 LLM agent 不是各自独立决策，而是通过多轮自然语言对话互相协商（一个提出计划、其他检查冲突、有冲突就“辩论”，直到达成无冲突的联合计划），这本质上是用“语言对话”实现了 \cref{sec:mh-why} 说的“难建模的高层协调决策”，只不过协调机制从 MARL 的奖励优化换成了 LLM 的对话推理；\emph{控制层（MPC）}——和 Level 2 完全一样。

\textbf{符号-几何接地（grounding）问题}：LLM 输出的是符号子目标（“抓取红块”），但 MPC 需要几何目标（“末端执行器移动到坐标 $(0.3,0.2,0.1)$”）。这中间需要一个\emph{接地}步骤，把符号映射到几何——通常由感知模块（物体检测 $+$ 位姿估计）完成。这个接地接口的可靠性是 LLM 混合系统的一个关键工程点（也是新的失败来源）。

\begin{insight}[协调知识的来源：从经验数据到人类知识]
RoCo 的“辩证协商”和 MARL 的“集中训练”是解决同一个问题（多机协调）的两种范式，根本差异在\textbf{协调知识的来源}。MARL 的协调知识来自\emph{海量交互数据 $+$ 奖励信号}（从经验中学）；LLM 的协调知识来自\emph{预训练语料 $+$ 常识推理}（从人类知识中迁移）。前者需要为每个任务训练、但能精确优化；后者零样本就能用、能理解语义、但不保证最优。这揭示了一个深刻趋势——混合范式的“高层决策器”正在从“任务专用的学习模型”走向“通用的基础模型”，从“为每个任务训练”走向“用一个通用模型零样本处理多任务”。
\end{insight}

\subsection{LLM 高层 vs MARL 高层：何时用哪个；新机遇与新风险}
\label{sec:mh-llm-vs-marl}

\begin{table}[H]\centering\small
\caption{LLM 高层 vs MARL 高层选型对比}
\label{tab:mh-llm-vs-marl}
\begin{tabular}{p{3.0cm}p{4.4cm}p{4.6cm}}
\toprule
维度 & MARL 高层 & LLM 高层（RoCo） \\
\midrule
协调知识来源 & 交互数据 $+$ 奖励 & 预训练常识 $+$ 推理 \\
是否需训练 & 是（每任务训练/微调） & 否（零样本/少样本） \\
决策频率 & 快（$\sim 10$ Hz 可达） & 慢（秒级，LLM 推理延迟） \\
常识利用 & 弱（需编码进奖励） & 强（内置海量常识） \\
自然语言指令 & 难 & 易（原生支持） \\
决策可解释 & 弱（黑盒） & 较强（思维链/对话可读） \\
最优性保证 & 较强（优化奖励） & 弱（推理非最优） \\
数值精度 & 好 & 差（LLM 弱项） \\
适用任务 & 高频协调、需精确优化、任务固定 & 低频符号决策、需常识/语义理解、任务开放多变 \\
\bottomrule
\end{tabular}
\end{table}

\noindent\textbf{选型口诀}：需要\emph{高频精确协调}（如密集编队避障）$\to$ MARL 高层；需要\emph{理解开放语义指令 $+$ 利用常识做符号级决策}（如“帮我把桌子收拾干净”）$\to$ LLM 高层。两者甚至可以叠加——LLM 做最顶层的任务级符号决策（慢、开放），MARL 做中层的运动级协调（快、精确），MPC 做底层执行。LLM 高层带来新能力，但也引入本章前面没出现过的\textbf{新风险}（\cref{tab:mh-llm-risk}）。

\begin{table}[H]\centering\small
\caption{LLM 高层的新风险与缓解思路}
\label{tab:mh-llm-risk}
\begin{tabular}{p{2.6cm}p{4.4cm}p{4.6cm}}
\toprule
新风险 & 描述 & 缓解思路 \\
\midrule
幻觉 & LLM 可能自信地输出错误/不存在的指令（如抓取场景里没有的物体） & 接地步骤做合法性检查；底层 MPC 对不可达目标降级；CBF 兜底物理安全 \\
延迟 & LLM 推理秒级，无法高频决策 & 只用于最高层低频符号决策；中层用 MARL/规则补高频协调 \\
接地脆弱性 & 符号$\to$几何接地依赖感知，感知错误会让正确的符号决策执行到错误几何目标 & 接地接口加置信度阈值；感知不确定时请求 LLM 重新决策 \\
不可证明 & LLM 决策完全不可形式化验证 & LLM 只做“建议”，所有安全约束仍由 CBF/MPC 硬保证——LLM 永远不碰安全层 \\
\bottomrule
\end{tabular}
\end{table}

\begin{insight}[一条铁律：安全永远不能交给高层决策器]
无论高层决策器多“聪明”（从 MARL 到 LLM），本章的一条铁律始终不变——\textbf{安全永远不能交给高层决策器，必须由底层的 CBF/MPC 硬保证}。MARL 会输出 OOD 危险动作，LLM 会幻觉出错误指令，它们的失效模式不同，但应对之策相同：让它们只做“建议性”的高层决策，把“绝对不可违反的物理安全”牢牢锁在底层的模型化安全层里。这呼应了 \cref{sec:mh-safety} 的核心——安全过滤层是横切所有控制路径的守门员，它不关心上面是策略网络还是 LLM，只保证任何下发到执行器的控制都不越界。这就是混合范式的安全哲学：\textbf{高层可以任意聪明也任意会犯错，底层保证它的错误永远不会变成物理灾难。}
\end{insight}

\begin{pitfall}[思维陷阱：以为 LLM 强大就能让它直接控制机器人底层]
错误做法：让 LLM 直接输出关节力矩、速度指令或高频控制信号。后果：频率完全不匹配（LLM 秒级 vs 控制毫秒级），LLM 无动力学概念输出物理上荒谬的控制，幻觉导致危险动作，系统无法运行或失控。根本原因：和“RL 不应出底层动作”同理但更严重——LLM 是符号/语言推理器，不是连续控制器。正确做法：LLM 只做最高层的低频符号决策（任务分解、协商），把符号子目标经接地变成几何目标交给底层 MPC，LLM 永远不碰底层。
\end{pitfall}

\begin{pitfall}[概念误区：把 LLM 的决策当成可信的，不做合法性检查和安全兜底]
错误做法：直接执行 LLM 输出的子目标，认为“LLM 这么聪明应该不会错”。后果：LLM 幻觉出一个不存在的物体/不可达的目标，系统盲目执行导致行为异常甚至危险。根本原因：LLM 会幻觉，且它的输出完全不可形式化验证。正确做法：LLM 输出只当“建议”——接地步骤做合法性检查（目标是否存在、是否可达），底层 MPC 对不可达目标降级，CBF 兜底物理安全。安全永远不依赖 LLM 的正确性。
\end{pitfall}

\begin{practice}
\begin{enumerate}
  \item \textbf{【架构分析】} 把 RoCo 的三层架构对齐到 \cref{sec:mh-hierarchy} 的四层架构上：RoCo 的哪几层对应 Level 3？控制层对应哪一层？RoCo 缺少了四层架构的哪一层（提示：底层执行频率）？如果要把 RoCo 部署到真实力控机器人上，你需要补上什么？
  \item \textbf{【对比辨析】} 针对两个任务分别判断用 MARL 高层还是 LLM 高层，并沿 \cref{tab:mh-llm-vs-marl} 的维度说明理由：(a) $50$ 架无人机表演密集队形变换（需高频、精确、可重复的协调）；(b) $3$ 台家用机器人协作“准备一顿晚餐”（开放任务、需常识理解“先洗菜再切菜”、需理解主人的自然语言偏好）。
  \item \textbf{【安全设计 $+$ 跨节综合】} 综合 \cref{sec:mh-safety} 与本节，为一个“LLM 高层 $+$ MPC 底层”的双机械臂协作系统设计完整的安全方案。回答：(a) LLM 幻觉出一个会导致两臂碰撞的子目标时，谁来阻止碰撞？(b) 接地步骤如何检查 LLM 子目标的合法性？(c) “安全绝不能交给 LLM”这条铁律在这个系统里如何具体落实？画出数据流，标出每一道安全检查的位置。
\end{enumerate}
\end{practice}

% ════════════════════════════════════════════════════════════════
\section*{本章常见误解汇总}
\addcontentsline{toc}{section}{本章常见误解汇总}

下表汇总本章贯穿始终、最容易绊倒初学者的认知误区。每一条都对应正文里详细论证过的本质洞察——如果某条你还觉得“不就是这样吗”，建议回对应小节重读。

\begin{table}[H]\centering\footnotesize
\caption{MARL 与传统规控混合常见误解辨正}
\label{tab:mh-misconceptions}
\begin{tabular}{p{4.4cm}p{6.4cm}p{1.8cm}}
\toprule
误解 & 正确理解 & 对应小节 \\
\midrule
把碰撞惩罚调到无穷大就能让 MARL 安全 & 奖励是期望意义的软目标，碰撞是它优化的一个“项”；安全必须是优化的“约束”（逐状态硬约束或期望约束），是数学层次的不同，不是调参能解决的 & \cref{sec:mh-why} \\
纯 MPC 更安全，所以应尽量多用 MPC & 安全和决策质量是正交维度。MPC 在“有模型执行”上安全，但在“难建模协调决策”上毫无优势——它只是把“学策略”换成“人脑手工设计代价” & \cref{sec:mh-why} \\
混合就是“先跑 MARL 再跑 MPC”的串联 & 混合的核心是“沿正确维度切分 $+$ 设计语义清晰的窄接口”，难点 $90\%$ 在接口设计。简单串联会导致语义不对齐、时间尺度错配 & \cref{sec:mh-why,sec:mh-interface} \\
让相邻层频率接近能提高响应速度 & 分层解耦依赖时间尺度分离（内环比外环快 $5$--$10$ 倍）。频率太接近会让两层强耦合、产生振荡 & \cref{sec:mh-hierarchy} \\
RL 应直接输出底层动作（力矩/轨迹） & RL 应只出低维运动学意图，把“展开成可行控制”留给有模型的 MPC。直接出底层动作既维度爆炸又让 RL 独自承受 sim-to-real gap & \cref{sec:mh-interface} \\
训练和部署用不同命令语义没关系 & MARL 是在某种“命令$\to$行为”映射下学的，这个映射由 MPC 如何解释命令决定。换了解释 $=$ 换了环境，策略失效 & \cref{sec:mh-interface} \\
对二阶/三阶系统用一阶 CBF & 当 $h$ 是位置函数、$u$ 是加速度/力时相对度 $\ge2$，$\dot h$ 不含 $u$，一阶 CBF 形同虚设，必须用 HOCBF & \cref{sec:mh-safety} \\
CBF-QP infeasible 时退回名义控制 & infeasible 恰恰发生在最危险时刻，退回名义控制 $=$ 在最该保护时撤掉保护。应降级到保守安全动作（急停/最大减速） & \cref{sec:mh-safety} \\
有了安全过滤就让上层完全不管安全 & 安全过滤保证“不越界”但不保证“行为优雅”。上层完全不管会让过滤器满负荷工作、系统在安全边缘横跳。上层仍需温和的安全意识 & \cref{sec:mh-safety} \\
CBF 保证安全就等于保证任务能完成 & CBF 保证安全（前向不变性）但不保证活性（liveness），二者是独立性质，可能死锁，需上层决策打破对称 & \cref{sec:mh-safety} \\
残差策略随便初始化就行 & 残差必须初始接近零（末层权重 $\approx0$），系统才从“纯基线行为”出发。随机初始化第一步就破坏基线 & \cref{sec:mh-residual} \\
残差可以无界 & 残差 RL 的安全性来自“不偏离可信基线太远”，必须用 \texttt{tanh}$\times\rho$ 限幅，否则退化成不安全的纯 RL & \cref{sec:mh-residual} \\
基线控制器也可以一起训练 & 基线的价值在于“可信、固定、有保证”，是残差 RL 的安全锚点。让它可训练 $=$ 锚点漂移 $=$ 安全论证失效 & \cref{sec:mh-residual} \\
Safe MARL 的期望安全 $=$ 硬安全保证 & 期望约束 $\mathbb{E}[\text{cost}]\le\delta$ 只保证平均安全，单条轨迹仍可能违反。安全关键场景必须再叠 CBF 硬过滤 & \cref{sec:mh-safemarl} \\
混合一定比纯方案好 & 混合的代价是工程复杂度。只有任务确实同时需要多种范式优势时混合才值得。简单任务用纯方案更好 & \cref{sec:mh-safemarl} \\
LLM 强大就能让它直接控制机器人底层 & LLM 是符号/语言推理器，频率（秒级）、动力学理解、数值精度全面不适合底层控制。它只能做最高层低频符号决策 & \cref{sec:mh-llm} \\
LLM 的决策可信，无需安全兜底 & LLM 会幻觉且不可形式化验证。它的输出只当“建议”，须做合法性接地检查 $+$ CBF 兜底物理安全 & \cref{sec:mh-llm} \\
\bottomrule
\end{tabular}
\end{table}

% ════════════════════════════════════════════════════════════════
\section*{本章小结}
\addcontentsline{toc}{section}{本章小结}

本章把整个多机器人方向的两条技术路线——学习（MARL）与模型（MPC/CBF/MAPF）——焊接成一个可部署的混合系统，并以此收口整个方向。核心逻辑链条是：
\begin{enumerate}
  \item \textbf{为什么混合}（\cref{sec:mh-why}）：纯 MARL 无硬安全保证、sim-to-real gap 弥漫、黑盒难调试；纯规控写不出难建模的高层协调决策。两者各有不可在单一范式内调和的天花板——这就是“探索-保证两难”。化解之道是\textbf{解耦原理}：沿“学习/模型”维度切分，让每种范式只在自己有比较优势的地方工作，通过窄接口通信，把 sim-to-real gap \emph{局部化}到高层决策接口上。
  \item \textbf{分层混合架构}（\cref{sec:mh-hierarchy}）：四层结构（MARL 高层 $10$ Hz $\to$ MPC 中层 $50$ Hz $\to$ WBC 底层 $200$--$1000$ Hz $+$ CBF 横切安全层），频率阶梯由\textbf{时间尺度分离}原理决定，每层可独立降级把单点故障隔离在层内。
  \item \textbf{RL$\to$MPC 接口}（\cref{sec:mh-interface}）：混合成败 $90\%$ 取决于接口设计。命令空间应满足低维、运动学抽象、可执行三原则。三种典型接口：目标速度、编队偏移、子目标。进阶接口：RL 出 MPC 尾代价（补“看不远”）。
  \item \textbf{安全过滤}（\cref{sec:mh-safety}）：CBF-QP 把安全做成可叠加在任何控制器上的最小修正（正交投影到安全半空间，\cref{thm:mh-cbf-closed}），提供逐状态硬保证（前向不变性，\cref{thm:mh-fwd-invariance}）；GCBF+ 用 GNN 参数化 CBF 扩展到上千 agent。
  \item \textbf{残差 RL}（\cref{sec:mh-residual}）：在可信基线上只学修正 $u=u_{\text{base}}+u_{\text{res}}^{\theta}$，样本效率高、有安全性下界（\cref{prop:mh-residual-bound}：残差限幅 $+$ 基线裕度）。多机场景里残差恰好捕捉“难建模的多体交互”。
  \item \textbf{Safe MARL 与选型}（\cref{sec:mh-safemarl}）：CMG $+$ Lagrangian 把安全写进训练目标（期望约束，$\mu$ 自适应权衡，\cref{thm:mh-lagrangian}），与 CBF 过滤（逐状态硬保证）互补成两道防线。三范式对比表 $+$ 四问选型框架收口整个方向。
  \item \textbf{前沿延伸——LLM 高层}（\cref{sec:mh-llm}）：高层决策器从 MARL 策略网络延伸到 LLM（RoCo 范式），底层仍是 MPC/CBF。这印证了四层架构的通用性——Level 3 是可替换插件。但无论高层多聪明，\textbf{安全永远由底层硬保证}，绝不交给会幻觉的 LLM。
\end{enumerate}

\noindent 贯穿全章的元洞察：\textbf{没有银弹，只有组合}。混合范式教的不是某个算法，而是“如何判断工具的比较优势、沿正确维度切分问题、设计模块接口”的系统设计方法论。从 MARL 高层到 LLM 高层的延伸进一步揭示：\emph{这个方法论的骨架是稳定的，可替换的只是每一层的具体实现}——这正是一个好的系统设计该有的样子。

\subsection*{术语速查表}
\begin{table}[H]\centering\small
\caption{本章核心术语（中/英）一句话定义}
\label{tab:mh-glossary}
\begin{tabular}{p{4.8cm}p{8.4cm}}
\toprule
术语（中/英） & 一句话定义 \\
\midrule
混合范式（Hybrid Paradigm） & 把 MARL 的探索能力与传统规控的可证明保证组合在一个系统里 \\
探索-保证两难（Exploration-Guarantee Dilemma） & RL 善探索无保证、控制有保证难探索，单一范式无法兼得 \\
解耦原理（Decoupling Principle） & 沿特征维度切分系统，每层只处理有比较优势的子问题，窄接口通信 \\
时间尺度分离（Time-Scale Separation） & 内环比外环快 $5$--$10$ 倍，外环可把内环近似当瞬时执行器，从而解耦 \\
命令空间（Command Space） & RL 输出给底层的语义量空间，应低维、运动学抽象、可执行 \\
安全过滤（Safety Filtering / Shielding） & 在名义控制执行前做最小修正，投影到安全集 \\
控制屏障函数（CBF） & 用 $h(x)\ge0$ 定义安全集，逐时刻约束 $\dot h\ge-\alpha(h)$ 保证前向不变 \\
前向不变性（Forward Invariance） & 初始在安全集内且每时刻满足 CBF 条件 $\Rightarrow$ 永远在安全集内 \\
相对度（Relative Degree） & $h$ 对 $u$ 求几阶导才显含 $u$；$\ge2$ 时需高阶 CBF \\
高阶 CBF（HOCBF） & 相对度 $\ge2$ 时构造嵌套函数链逼出对 $u$ 的约束 \\
CBF-QP & 求“最接近名义控制且满足 CBF 约束”的控制的二次规划 \\
GCBF+（Graph CBF+） & 用 GNN 参数化 CBF，可扩展到上千 agent 的分布式安全过滤 \\
残差强化学习（Residual RL） & 在固定基线控制器上叠加 RL 学的修正 $u=u_{\text{base}}+u_{\text{res}}$ \\
安全性下界（Safety Lower Bound） & 残差限幅 $+$ 基线裕度保证总控制不偏离安全边界 \\
约束马尔可夫博弈（CMG） & 多智能体的约束优化：max 奖励 s.t.\ 期望累积代价 $\le$ 阈值 \\
Lagrangian 松弛 & 用乘子 $\mu$ 把约束吸进目标，原始-对偶交替优化 \\
原始-对偶优化（Primal-Dual） & 交替更新策略 $\theta$（原始）与乘子 $\mu$（对偶），$\mu$ 应更新更慢 \\
符号-几何接地（grounding） & LLM 符号子目标 $\to$ MPC 几何目标的映射，LLM 混合系统的新失败来源 \\
\bottomrule
\end{tabular}
\end{table}

\subsection*{知识点总表}
\begin{table}[H]\centering\footnotesize
\caption{本章知识点与难度}
\label{tab:mh-knowledge}
\begin{tabular}{cp{3.6cm}p{6.0cm}cc}
\toprule
编号 & 知识点 & 核心要点 & 节 & 难度 \\
\midrule
13-1 & 探索-保证两难 & 纯 MARL 无硬保证、纯规控难建模协调，单范式有天花板 & \cref{sec:mh-why} & $\star\star\star$ \\
13-2 & 解耦原理 & 沿比较优势维度切分 $+$ 窄接口，gap 局部化 & \cref{sec:mh-why} & $\star\star\star$ \\
13-3 & 四类 sim-to-real gap & 动力学/感知/执行/多体交互，各甩给最适合的层 & \cref{sec:mh-why} & $\star\star$ \\
13-4 & 四层混合架构 & MARL $10$Hz $\to$ MPC $50$Hz $\to$ WBC $200$--$1000$Hz $+$ CBF 横切 & \cref{sec:mh-hierarchy} & $\star\star\star$ \\
13-5 & 时间尺度分离 & 相邻层差一个数量级，外环把内环当瞬时执行器 & \cref{sec:mh-hierarchy} & $\star\star\star$ \\
13-6 & 分层降级 & 每层在上层失效时优雅降级，单点故障隔离在层内 & \cref{sec:mh-hierarchy} & $\star\star$ \\
13-7 & 命令空间三原则 & 低维、运动学抽象、可执行 & \cref{sec:mh-interface} & $\star\star\star\star$ \\
13-8 & 三种 RL$\to$MPC 接口 & 目标速度 / 编队偏移 / 子目标 & \cref{sec:mh-interface} & $\star\star\star$ \\
13-9 & RL 出 MPC 尾代价 & RL 学 Q 函数补 MPC“看不远”的短板 & \cref{sec:mh-interface} & $\star\star\star\star$ \\
13-10 & CBF 与前向不变性 & $\dot h\ge-\alpha(h)$ 保证安全集前向不变 & \cref{sec:mh-safety} & $\star\star\star$ \\
13-11 & CBF-QP 闭式解 & $u_{\text{safe}}=u_{\text{nom}}-\frac{\max(0,a^\top u_{\text{nom}}-b)}{\|a\|^2}a$，正交投影 & \cref{sec:mh-safety} & $\star\star\star\star$ \\
13-12 & HOCBF & 相对度 $\ge2$ 时的嵌套函数链构造 & \cref{sec:mh-safety} & $\star\star\star\star$ \\
13-13 & GCBF+ & GNN 参数化 CBF，置换等变 $+$ 邻居数不变，扩展到上千 agent & \cref{sec:mh-safety} & $\star\star\star\star$ \\
13-14 & 残差 RL 公式 & $u=u_{\text{base}}+u_{\text{res}}^{\theta}$，学修正而非完整策略 & \cref{sec:mh-residual} & $\star\star\star\star$ \\
13-15 & 残差初始化与限幅 & 末层权重 $\approx0$ $+$ \texttt{tanh}$\times\rho$ 限幅，保证从基线出发且有界 & \cref{sec:mh-residual} & $\star\star\star$ \\
13-16 & 残差安全性下界 & 残差有界 $+$ 基线裕度 $\Rightarrow$ 全程不越界 & \cref{sec:mh-residual} & $\star\star\star$ \\
13-17 & 约束马尔可夫博弈 & max 奖励 s.t.\ $\mathbb{E}[\sum\gamma^tc]\le\delta$ & \cref{sec:mh-safemarl} & $\star\star\star$ \\
13-18 & Lagrangian 原始-对偶 & $\theta$ 升 $A^r-\mu A^c$，$\mu$ 升约束违反量，$\mu$ 更新更慢 & \cref{sec:mh-safemarl} & $\star\star\star\star$ \\
13-19 & 两种安全语义 & 期望安全（训练内化，软）vs 逐状态硬保证（部署过滤，硬） & \cref{sec:mh-safemarl} & $\star\star\star$ \\
13-20 & 三范式选型框架 & 沿“硬安全/难建模决策/有无基线/规模”四问决策 & \cref{sec:mh-safemarl} & $\star\star\star$ \\
13-21 & LLM 高层（RoCo） & LLM 替换 MARL 做高层符号决策，辩证协商 $+$ MPC 底层 & \cref{sec:mh-llm} & $\star\star\star\star$ \\
13-22 & 高层决策器的可替换性 & Level 3 是可替换插件（策略网络/LLM），底层安全机制不变 & \cref{sec:mh-llm} & $\star\star\star$ \\
13-23 & 符号-几何接地 & LLM 符号子目标 $\to$ MPC 几何目标的映射，新失败来源 & \cref{sec:mh-llm} & $\star\star\star$ \\
13-24 & 安全永不交给高层 & 无论高层多聪明（MARL/LLM），物理安全必须由底层 CBF/MPC 硬保证 & \cref{sec:mh-safety,sec:mh-llm} & $\star\star\star$ \\
\bottomrule
\end{tabular}
\end{table}

% ════════════════════════════════════════════════════════════════
\section*{延伸阅读}
\addcontentsline{toc}{section}{延伸阅读}

\begin{itemize}
  \item \textbf{综述与理论地图}：Brunke, Greeff 等\emph{Safe Learning in Robotics: From Learning-Based Control to Safe Reinforcement Learning}（Annual Review of Control, 2022\cite{brunke2022safe}）——本章的理论地图，系统梳理安全 RL 的三条路线（安全过滤 / 约束优化 / 可达性）。
  \item \textbf{安全过滤一侧}：GCBF+\cite{zhang2025gcbfplus}（T-RO 2024，\cref{sec:mh-safety} 的核心论文，GNN 参数化 CBF 扩展到上千 agent，项目主页有代码和视频）；Ames 等 CBF 综述\cite{ames2019cbf}（ECC 2019，理论根基）；CBF-RL\cite{gandhi2023shield} 与 SHIELD\cite{yin2023shield}（2025，把安全从部署期推进到训练期内化 / 推广到学习动力学）。
  \item \textbf{分层混合一侧}：Bjelonic 等\emph{Combining MPC and Predictive RL for Quadruped Locomotion}\cite{bjelonic2023mpcrl}（2023，\cref{sec:mh-interface} “RL 出尾代价”接口的代表作）；RoCo\cite{mandi2024roco}（ICRA 2024，混合范式向“LLM 高层”延伸）。
  \item \textbf{残差学习一侧}：Johannink 等\emph{Residual Reinforcement Learning for Robot Control}\cite{johannink2019residual}（ICRA 2019，\cref{sec:mh-residual} 的奠基论文，必读）；Residual MPC\cite{residualmpc2025}（2025，把残差的基线升级为 GPU 并行 MPC）；ResNet\cite{he2016resnet}（残差思想的源头）。
  \item \textbf{约束 RL 理论}：Altman\emph{Constrained Markov Decision Processes}\cite{altman1999cmdp}（1999，CMDP 原典，\cref{sec:mh-safemarl} 的根基）；Achiam 等 CPO\cite{achiam2017constrained}（ICML 2017，信赖域内保证约束满足，Lagrangian 之外的另一条路）。
  \item \textbf{开源代码库}：\verb@MIT-REALM/gcbfplus@（GCBF+ 官方实现，JAX）；\verb@PKU-MARL/HARL@（HAPPO/MAPPO 等，含 safe MARL 的 Lagrangian 实现）；\verb@osqp/osqp@（CBF-QP 用的 QP 求解器）；\verb@leggedrobotics/ocs2@（MPC（SQP-RTI）$+$ WBC，工业级实现）。
\end{itemize}

\section*{本章与后续章节的关系}
\addcontentsline{toc}{section}{本章与后续章节的关系}

本章是多机器人方向的\textbf{收口章}。如果说前面各章是在“打造工具”（通信、共识、分配、规划、MPC、MARL、CBF 各是一件工具），本章就是教你“如何把工具组装成机器”——它不引入太多全新的算法，而是提供组装这些工具的\textbf{系统设计方法论}。
\begin{itemize}
  \item \textbf{承接 \cref{ch:plan-safemarl}}：那里给出 MARL 全栈（CTDE/非平稳/HAPPO/MAPPO/值分解）与 CBF$+$博弈的安全证书；本章把这些作为高层决策器与安全过滤层的现成零件，组装进分层混合系统。
  \item \textbf{承接 \cref{ch:mr-consensus}}：共识与 ADMM 是分布式 MPC（本章 Level 2 的分布式版）与多机协调的底层机制。
  \item \textbf{延伸到方向专题}：本章的安全过滤（\cref{sec:mh-safety}）、残差（\cref{sec:mh-residual}）是方向无关的通用组件，无人机/足式/机械臂专题会针对各自动力学特化它们（如足式用单刚体动力学 CBF、无人机用几何控制残差）。
\end{itemize}
读完本章，你应完成从“会用单个算法”到“会设计多机器人系统”的认知跃迁。

% ════════════════════════════════════════════════════════════════
\section*{故障排查手册}
\addcontentsline{toc}{section}{故障排查手册}

混合系统的调试难点在于\textbf{故障可能源自任何一层或层间接口}，所以排查的第一原则是“逐层隔离”——先确定故障在哪一层，再深入。

\begin{description}
  \item[故障 1：仿真训练完美，部署真机频繁碰撞。] \emph{可能原因}：(a) 没有部署安全过滤层，纯靠 MARL 软安全；(b) CBF 相对度搞错（二阶系统用了一阶 CBF，过滤形同虚设）；(c) CBF 的 $\alpha$ 太大，贴边才刹来不及；(d) sim-to-real gap 让 MARL 输出 OOD 危险动作而无过滤兜底。\emph{排查}：1. 确认安全过滤层是否真在部署回路里（打印每帧是否调用、修正量是否非零）；2. 检查 CBF 相对度——控制是力/加速度且 $h$ 是位置函数则必须 HOCBF；3. 注入一个故意冲向障碍的名义控制，验证过滤器是否拦截、何时开始减速；4. 减小 $\alpha$ 让减速更早。\emph{相关}：\cref{sec:mh-why,sec:mh-safety}。
  \item[故障 2：机器人整体协调但每台高频抖动。] \emph{可能原因}：(a) 相邻层频率太接近，时间尺度分离失效；(b) RL$\to$MPC 命令未限幅/未平滑，命令跳变导致 MPC 输出剧烈饱和；(c) 训练与部署的接口语义不一致；(d) CBF 过滤与名义控制持续对抗。\emph{排查}：1. 检查各层频率比是否 $\ge5$ 倍；2. 打印 RL 输出的命令序列看是否跳变、确认限幅和速率限制生效；3. 核对训练时和部署时的接口代码是否\emph{逐行一致}；4. 打印 CBF 修正量，若持续大幅修正说明上层太激进。\emph{相关}：\cref{sec:mh-hierarchy,sec:mh-interface,sec:mh-safety}。
  \item[故障 3：CBF-QP 频繁求解失败（infeasible）。] \emph{可能原因}：(a) 多个 pairwise 约束在密集场景下互相冲突；(b) $\alpha$/安全半径过保守把可行域压没了；(c) 没有 infeasible 降级逻辑；(d) 缺少打破对称的上层决策导致死锁。\emph{排查}：1. 统计 infeasible 发生时的邻居数和密度；2. 检查是否有 infeasible 降级（应降到急停/最大减速，不是退回名义控制）；3. 考虑用软约束 CBF（加松弛变量 $+$ 大惩罚）保证 QP 始终可行；4. 检查上层是否提供“谁让谁”的协调决策打破对称。\emph{相关}：\cref{sec:mh-safety}。
  \item[故障 4：残差 RL 训练不收敛或行为比基线还差。] \emph{可能原因}：(a) 残差策略末层随机初始化，第一步就破坏基线；(b) 残差不限幅，偏离基线太远；(c) 基线被误加进训练，污染可信性；(d) 基线太弱，残差不“小”。\emph{排查}：1. 打印训练第一步的残差幅值，应接近零；2. 确认 \texttt{tanh}$\times\rho$ 限幅生效；3. 检查优化器参数列表，确认基线参数\emph{未}被包含；4. 评估基线单独运行的性能，太差则先用控制论改进基线。\emph{相关}：\cref{sec:mh-residual}。
  \item[故障 5：Safe MARL 训练中安全代价振荡、约束时满足时违反。] \emph{可能原因}：(a) Lagrangian 乘子 $\mu$ 学习率太大，破坏原始-对偶的时间尺度分离；(b) 安全代价 $c_{i,t}$ 设计不合理（量纲/尺度问题）；(c) 把期望安全误当硬保证，没有运行时过滤兜底。\emph{排查}：1. 把 $\mu$ 的学习率降到策略学习率的 $1/10$--$1/100$；2. 打印 $\mu$ 的轨迹，应缓慢单调趋稳而非震荡；3. 检查各代价分量的尺度是否可比；4. 评估时叠加 CBF 过滤做硬兜底。\emph{相关}：\cref{sec:mh-safemarl}。
  \item[故障 6（前沿）：LLM 高层系统执行了错误/危险的目标。] \emph{可能原因}：(a) LLM 幻觉出错误指令且未做合法性检查；(b) 符号$\to$几何接地错误；(c) 误把 LLM 输出当可信、未做安全兜底；(d) LLM 直接碰了不该碰的层。\emph{排查}：1. 在接地步骤打印“LLM 符号目标 $\to$ 几何目标”的映射，检查目标是否存在、是否可达；2. 核对感知模块对该物体的检测/位姿是否正确；3. 确认 CBF 过滤层在 LLM 系统里仍生效；4. 确认 LLM 只输出高层符号决策、未触碰底层控制。\emph{相关}：\cref{sec:mh-llm,sec:mh-safety}。
\end{description}
